% =============================================================
% Chapter 9  C 端应用（四）：AI 伴侣、搜索与个人助手
% Book: The AI Startup Landscape
% =============================================================

\chapter{C 端应用（四）：AI 伴侣、搜索与个人助手}
\label{ch:consumer-social}

% AI as companion, search engine, and OS-level intelligence.
% The most intimate and pervasive forms of consumer AI.

如果说前几章讨论的AI工具——办公、创作、教育——都在帮用户\textbf{做事}，
那么本章涉及的AI形态则更加根本：它们试图成为用户的\textbf{对话伙伴}、
\textbf{信息入口}和\textbf{系统级智能}。从ChatGPT到豆包，从
Character.ai到Perplexity，从Apple Intelligence到Windows Copilot，
这些产品覆盖了人与AI交互最高频、最私密、最具黏性的场景。

这也是整个AI应用层中商业模式最多样、争议最大、格局变化最快的领域。
一个ChatGPT级别的通用助手需要每年数十亿美元的基础设施投入，一个
AI伴侣应用却可能只用几百万美元就触达千万级用户。搜索引擎的颠覆者
正在挑战Google二十年来的广告霸权，而操作系统级AI的嵌入则意味着——
未来你甚至不需要打开任何App，智能就已经在那里了。

本章还延伸至三个与"AI交互"密切相关的前台业务场景——销售、客服和
HR招聘。这些领域的共同特征是：大量重复性的人际沟通、丰富的交互数据、
以及对效率提升的强烈需求。AI在这些场景中不是替代创造力，而是在
\textbf{替代重复劳动}，并正在从辅助工具进化为自主运行的AI代理。

从市场规模的角度看，本章涵盖的赛道合计代表了AI应用层中最大的
价值池。仅以几个子赛道为例：

\begin{itemize}
  \item \textbf{AI聊天机器人/助手}：ChatGPT+Claude+Gemini等的合计
    年化收入已超过200亿美元，且保持10倍级的年增长
  \item \textbf{AI伴侣}：2024年市场规模108亿美元，预计2034年达到
    2908亿美元（CAGR 39\%）
  \item \textbf{AI搜索}：Google的搜索广告收入每年超过1700亿美元
    ——AI搜索即使只争取其中5\%的份额，也是一个85亿美元的市场
  \item \textbf{AI客服}：2026年市场约100--110亿美元，CAGR 23--26\%
  \item \textbf{收入智能/销售AI}：Gong ARR已超过3亿美元，Clay
    估值31亿美元，整个赛道年增速超过40\%
  \item \textbf{AI HR}：Eightfold估值21亿美元，Paradox被Workday
    收购，市场在快速整合
\end{itemize}

这些赛道虽然看似分散，但它们共享一个底层逻辑：\textbf{AI正在重新
定义人机交互的每一个入口}。从你早上醒来对手机说的第一句话
（Siri/Gemini），到工作中使用的销售工具（Gong/Clay），到晚上睡前
与AI伴侣的对话（Character.ai/Replika），AI正在渗透到生活和工作的
每一个环节。

\begin{corebox}[AI助手格局：从ChatGPT到系统级AI]
截至2026年3月，全球AI助手生态呈现出三层结构：
\begin{enumerate}
  \item \textbf{独立AI应用}：ChatGPT（周活9亿）、Claude（月活1890万）、
    豆包（月活2.27亿）、Gemini App、Kimi Chat等。用户主动打开一个
    独立应用进行对话。这是当前最成熟的交互模式。
  \item \textbf{嵌入式AI}：搜索引擎中的AI答案（Perplexity、秘塔）、
    浏览器内置AI（Brave Leo、Edge Copilot）、应用内AI功能
    （Notion AI、Canva AI）。AI作为已有产品的增强层存在。
  \item \textbf{系统级AI}：Apple Intelligence、Google Gemini Android、
    Samsung Galaxy AI、Windows Copilot。AI直接集成在操作系统中，
    无需额外安装，覆盖数十亿设备。
\end{enumerate}
这三层并非替代关系，而是共存与竞争。独立应用拥有最深的功能和最强的
品牌认知；嵌入式AI拥有场景优势；系统级AI拥有分发优势。最终的赢家
很可能是那些能同时覆盖多个层次的玩家——这也是为什么OpenAI要做
ChatGPT桌面客户端，Google要把Gemini从App做进Android内核，
而苹果则选择直接从系统层切入。
\end{corebox}

本章将依次讨论以下九大子赛道：

\begin{enumerate}
  \item \textbf{AI聊天机器人与通用助手}（\S\ref{sec:social-assistants}）：
    ChatGPT、Claude、Gemini、Pi/Inflection及中国助手（豆包、Kimi、千问等）
  \item \textbf{AI伴侣与社交}（\S\ref{sec:social-companions}）：
    Character.ai、Replika、Nomi、星野/Talkie、猫箱/Glow
  \item \textbf{AI角色扮演与互动小说}（\S\ref{sec:social-roleplay}）：
    NovelAI、JanitorAI、SillyTavern生态
  \item \textbf{AI搜索引擎}（\S\ref{sec:social-search}）：
    Perplexity、You.com、Kagi、Phind、秘塔AI搜索、天工AI
  \item \textbf{AI浏览器}（\S\ref{sec:social-browsers}）：
    Arc/Dia、Brave Leo、Opera AI、Microsoft Edge Copilot
  \item \textbf{OS级AI集成}（\S\ref{sec:social-os}）：
    Apple Intelligence、Google Gemini Android、Samsung Galaxy AI、Windows Copilot
  \item \textbf{AI销售与营销工具}（\S\ref{sec:social-sales}）：
    Gong、Clay、Apollo AI、Outreach、Chorus
  \item \textbf{AI客服}（\S\ref{sec:social-customer-service}）：
    Sierra AI、Intercom Fin、Zendesk AI、Ada、Forethought
  \item \textbf{AI HR与招聘}（\S\ref{sec:social-hr}）：
    HireVue、Paradox、Eightfold AI、Moonhub
\end{enumerate}

每个赛道我们将分析核心公司的融资、用户数据、产品差异化和竞争
格局。其中Character.ai和Perplexity作为AI伴侣和AI搜索的标杆公司，
将获得最详细的分析。中国市场的AI助手（特别是豆包和Kimi）和AI搜索
（秘塔、天工）也将得到深入覆盖。


% =============================================================
\section{AI 聊天机器人与通用助手}
\label{sec:social-assistants}
% =============================================================

通用AI助手是本轮AI浪潮中用户规模最大、增长最快的品类。这些产品的
核心功能看似简单——你问问题，AI回答——但在这个简单的交互背后，是
数百亿美元的基础设施投入、数十万张GPU的算力消耗、以及对人类语言和
知识最深层的理解。

本节梳理西方和中国市场最重要的通用AI助手，分析其用户规模、商业模式
和竞争策略。图~\ref{fig:assistant-landscape}展示了当前AI助手与伴侣产品
沿"通用性"和"情感深度"两个维度的分布格局。

\begin{figure}[htbp]
\centering
\begin{tikzpicture}[
  company/.style={
    font=\footnotesize, rounded corners=2pt,
    inner sep=3pt, minimum height=0.55cm,
  },
  axislbl/.style={
    font=\small\bfseries, text=figGray,
  },
  quadlbl/.style={
    font=\footnotesize\itshape, text=figGray!80,
  },
]

% --- Axes ---
\draw[figGray, line width=1pt, -{Stealth[length=5pt]}]
  (0,0) -- (12,0) node[axislbl, below left=2pt] {Generality};
\draw[figGray, line width=1pt, -{Stealth[length=5pt]}]
  (0,0) -- (0,8.5) node[axislbl, above left=2pt, rotate=90, anchor=south] {Emotional Depth};

% --- Axis endpoint labels ---
\node[font=\scriptsize, text=figGray] at (0.8,-0.4) {Vertical};
\node[font=\scriptsize, text=figGray] at (11.2,-0.4) {General-purpose};
\node[font=\scriptsize, text=figGray, rotate=90, anchor=south] at (-0.5,0.8) {Tool};
\node[font=\scriptsize, text=figGray, rotate=90, anchor=south] at (-0.5,7.5) {Companion};

% --- Quadrant dividers ---
\draw[figLightGray, line width=0.6pt, dashed] (6,0) -- (6,8.2);
\draw[figLightGray, line width=0.6pt, dashed] (0,4.2) -- (11.7,4.2);

% --- Quadrant labels ---
\node[quadlbl] at (3,7.8) {Niche Companion};
\node[quadlbl] at (9,7.8) {Social Companion};
\node[quadlbl] at (3,0.5) {Vertical Tool};
\node[quadlbl] at (9,0.5) {General Assistant};

% --- Bottom-left: Vertical tools ---
\node[company, fill=figBlue!12, draw=figBlue!50] at (2.5,2.0) {Perplexity};
\node[company, fill=figBlue!12, draw=figBlue!50] at (2.5,1.2) {Phind};

% --- Bottom-right: General assistants ---
\node[company, fill=figBlue!12, draw=figBlue!50] at (8.5,2.8) {ChatGPT};
\node[company, fill=figBlue!12, draw=figBlue!50] at (10.2,2.2) {Claude};
\node[company, fill=figBlue!12, draw=figBlue!50] at (8.5,1.5) {Gemini};
\node[company, fill=figGreen!12, draw=figGreen!50] at (10.2,1.2) {\kaishu 豆包};
\node[company, fill=figGreen!12, draw=figGreen!50] at (7.0,1.5) {Kimi};

% --- Top-left: Niche companion ---
\node[company, fill=figOrange!12, draw=figOrange!50] at (2.5,6.0) {NovelAI};
\node[company, fill=figOrange!12, draw=figOrange!50] at (2.5,5.2) {JanitorAI};

% --- Top-right: Social companion ---
\node[company, fill=figRed!12, draw=figRed!50] at (8.5,6.5) {Character.ai};
\node[company, fill=figRed!12, draw=figRed!50] at (8.5,5.5) {Replika};
\node[company, fill=figRed!12, draw=figRed!50] at (10.2,6.0) {\kaishu 星野};

% --- Market size annotations ---
\node[font=\tiny, text=figGray, align=center] at (9,3.6)
  {General assistants\\WAU $>$ 1B combined};
\node[font=\tiny, text=figGray, align=center] at (8.5,7.3)
  {AI companion market\\\$108B (2024) $\to$ \$2908B (2034)};

\end{tikzpicture}
\caption{AI助手与伴侣市场格局：通用性 vs 情感深度（2026年初）}
\label{fig:assistant-landscape}
\end{figure}

% -----------------------------------------------------------
\subsection{ChatGPT：定义品类的超级应用}
\label{subsec:chatgpt}
% -----------------------------------------------------------

2022年11月30日发布的ChatGPT不仅定义了"AI聊天机器人"这个品类，
更成为了科技史上增长最快的消费产品。从发布到1亿用户仅用了两个月
——相比之下，TikTok用了9个月，Instagram用了2.5年。

到2026年3月，ChatGPT的规模已经远超其他所有AI应用的总和。以下是
截至2026年初的关键数据：

\begin{itemize}
  \item \textbf{周活跃用户}：9亿（2026年2月27日OpenAI官方公布），
    从2025年2月的4亿和2025年10月的8亿持续攀升，距离10亿大关
    仅一步之遥
  \item \textbf{月度访问量}：57.2亿次（2026年1月），跻身全球
    访问量前五的网站
  \item \textbf{付费订阅用户}：5000万（2026年2月），其中
    ChatGPT Plus（\$20/月）约1000万，Team/Enterprise方案
    贡献了大量增量
  \item \textbf{年化经常性收入}：100亿美元（2026年初估计），
    较2024年底的约50亿美元翻倍
  \item \textbf{日均提示量}：25亿条
  \item \textbf{全球AI聊天机器人市场份额}：81.13\%
    （DemandSage, 2026年3月）
\end{itemize}

ChatGPT的商业模式经历了从纯订阅到多元化的演进。最初仅有\$20/月的
Plus方案，后来加入了\$200/月的Pro方案（面向重度用户和开发者、
提供o1 pro mode等高阶推理能力）、Team方案
（\$25/用户/月）和Enterprise方案（定制定价、SSO、
数据隔离等企业级功能）。2025年起，OpenAI还开始在免费版中测试广告
模式，并通过API收入实现B2B变现。根据TechCrunch的报道，OpenAI
曾考虑推出一个\$2000/月的超高端方案，面向研究机构和企业的最密集
使用场景。

OpenAI在2026年2月完成了1100亿美元融资（Amazon \$50B、NVIDIA \$30B、SoftBank \$30B），投后估值
达8400亿美元（\$840B）。值得注意的是，ChatGPT的消费者收入仅占OpenAI总收入的
一部分——API收入和企业合同同样是重要的营收来源。但ChatGPT的品牌
效应和用户基础是整个商业帝国的基石：它让"OpenAI"成为了AI的代名词，
就像Google曾经让自己的名字变成了"搜索"的动词。

ChatGPT的产品演进也值得关注。从最初的纯文本对话开始，ChatGPT在
两年多的时间里依次加入了：联网搜索（Browse with Bing）、代码
执行（Code Interpreter）、图像理解（GPT-4V）、图像生成
（DALL-E集成）、自定义GPTs（GPT Store）、语音对话（Advanced
Voice Mode）、视频理解、以及最新的Operator功能（AI自主操作
浏览器完成任务）。每一次功能扩展都在拓宽ChatGPT的使用场景，
从一个对话工具逐渐演化为一个\textbf{通用AI平台}。

2026年初，OpenAI还推出了ChatGPT的桌面客户端——支持macOS和Windows
——用户可以通过快捷键在任何应用中召唤ChatGPT。这一举措的战略
意图很明显：在苹果、Google和微软将AI嵌入操作系统之前，先占领
用户桌面的快捷入口。

OpenAI CEO Sam Altman表示，2026年1月和2月是ChatGPT历史上新增
付费订阅用户最多的两个月——这表明ChatGPT的增长不仅没有放缓，
反而在加速。"人们使用ChatGPT来学习、写作、规划和构建。随着使用
规模的扩大，产品在用户能立即感受到的方面不断改善：更快的响应、
更高的可靠性、更强的安全性和更一致的性能。"

ChatGPT面临的最大长期威胁不是任何单一竞品，而是\textbf{入口被
截流}。当Apple Intelligence让用户直接对Siri说"帮我写这封邮件"，
当Google Gemini在搜索结果中直接给出答案，当Microsoft Copilot
嵌入Office的每一个角落，用户访问ChatGPT的动机就被分流了。
OpenAI的应对策略是多管齐下：桌面客户端争夺快捷入口、语音模式争夺
语音入口、Operator争夺任务入口、与苹果合作争夺Siri溢出流量。
但这场"入口保卫战"的结果，将在很大程度上决定OpenAI未来的估值
是3000亿美元还是3万亿美元。

\begin{table}[htbp]
\centering
\caption{全球主要AI助手对比（截至2026年3月）}
\label{tab:ai-assistants-comparison}
\small
\begin{tabularx}{\textwidth}{l r r r X}
\toprule
\rowcolor{figLightGray}
\textbf{产品} & \textbf{用户规模} & \textbf{年化收入}
  & \textbf{公司估值} & \textbf{核心优势} \\
\midrule
ChatGPT     & 周活9亿
  & $\sim$\$10B & \$300B & 品牌=品类; 最全功能; 最大生态 \\
Claude      & 月活1890万
  & \$14B & \$380B & 安全优先; 代码; 企业; Claude Code \\
Gemini      & 未披露$^*$
  & 含于Google & \$2.3T$^\dagger$ & Google生态集成; 个人智能; Android \\
豆包        & 月活2.27亿
  & 未披露 & \$480B$^\dagger$ & 抖音生态; 多模态; 极致低价 \\
Kimi        & 月活$\sim$2500万
  & 快速增长 & \$18B & 长上下文; K2.5开源; 全球化 \\
DeepSeek    & 月活1.36亿
  & API为主 & 私有 & 开源模型; 推理优化; 行业推动 \\
Copilot     & 月活1.05亿
  & 含于MS & \$3.3T$^\dagger$ & Windows/Office集成; 企业渗透 \\
\bottomrule
\end{tabularx}

\medskip
\noindent {\scriptsize $^*$ Gemini作为Google搜索的组成部分，
用户规模难以独立衡量。$^\dagger$ 为母公司市值（Google/字节/微软），
非AI业务独立估值。数据来源：DemandSage, Sacra, SaaStr, Bloomberg,
QuestMobile, FirstPageSage. 截至2026年3月。}
\end{table}

% -----------------------------------------------------------
\subsection{Claude：安全优先的企业宠儿}
\label{subsec:claude-assistant}
% -----------------------------------------------------------

Anthropic的Claude走了一条与ChatGPT截然不同的路——\textbf{安全优先、
企业为重}。

Claude于2023年3月推出，在消费者市场的用户规模远小于ChatGPT：网页端
月活约1890万，移动端约738万（DemandSage, 2026年3月）。Claude在
AI聊天机器人市场的份额约为4.5\%。但如果只看这些消费端数字，你会
严重低估Anthropic的商业实力。

公司的真正战场在企业端——Anthropic报告有超过30万企业客户，年化收入
在2026年2月达到\textbf{140亿美元}，较2024年底的约10亿美元实现了
14倍增长。这一增速在企业软件史上前所未见：没有任何一家B2B软件公司
曾连续三年实现10倍以上的年增长。

这一惊人增速的核心驱动力之一是\textbf{Claude Code}——Anthropic于
2025年推出的编程和AI代理产品。Claude Code允许开发者通过命令行界面
直接与Claude交互来编写、调试和重构代码。到2026年初，Claude Code
已达到约25亿美元的年化收入。Claude在开发者群体中建立了极强的口碑，
被普遍认为在代码生成、长文档理解、指令遵从和"不自作主张"方面优于
竞品。

2026年2月，Anthropic完成了300亿美元的G轮融资——科技史上第二大
私募融资轮，仅次于OpenAI的1100亿美元——投后估值\textbf{3800亿美元}。
Google、Salesforce Ventures和Spark Capital等参与了投资。这一估值
意味着Anthropic的收入倍数约为27倍（基于140亿美元ARR），在SaaS公司中
处于极高水平，但考虑到其增速（10x/年），投资者显然认为这一估值
仍然合理。

Claude与ChatGPT的差异化可以从三个维度理解：

\textbf{第一，安全理念}。Anthropic由前OpenAI研究副总裁Dario
Amodei和姐姐Daniela Amodei创立，核心使命是开发"安全的AI"。
Claude的训练过程使用了"Constitutional AI"（宪法式AI）方法，
通过一组明确的原则来约束模型行为。这种方法论使得Claude在输出中
更少产生有害内容，也更倾向于在不确定时表示"我不知道"而非编造答案。

\textbf{第二，产品定位}。ChatGPT试图成为"万能工具"，功能边界
不断扩展；Claude则更专注于\textbf{专业工作场景}——长文档分析、
代码编写、研究辅助。Claude的200K token上下文窗口（约15万英文单词）
使其在处理长文档方面拥有显著优势。

\textbf{第三，商业策略}。OpenAI两条腿走路（消费者+企业），Anthropic
则将绝大部分资源押注在企业市场。Claude的API定价通常低于OpenAI的
同等模型，以低价策略争夺开发者和企业客户。

与ChatGPT的81\%市场份额相比，Claude的4.5\%看似微不足道。但如果
只看企业和开发者市场，格局完全不同——Claude在技术团队中的渗透率
远高于消费者市场，许多硅谷公司同时使用ChatGPT和Claude，后者
尤其被偏好用于编程和文档分析任务。

% -----------------------------------------------------------
\subsection{Google Gemini：搜索巨头的AI转身}
\label{subsec:gemini-assistant}
% -----------------------------------------------------------

Google的AI助手经历了一段曲折的品牌演进：从Google Assistant到Bard，
再到Gemini。2024年2月，Google将Bard正式更名为Gemini，并推出了基于
Gemini Ultra模型的付费方案Gemini Advanced（\$19.99/月，包含在
Google One AI Premium中）。

这次品牌重塑的背后是一次深刻的战略觉醒。2022年底ChatGPT的横空出世
让Google内部拉响了"红色警报"（code red）——这是自1998年成立以来，
Google第一次感受到搜索业务面临的生存级威胁。如果用户开始用AI助手
而不是搜索引擎来获取信息，Google每年超过1700亿美元的广告收入就
面临根本性挑战。

Gemini的独特优势在于其与Google生态系统的深度集成。2026年3月18日，
Google将Gemini的"个人智能"功能向\textbf{所有美国用户}开放——
此前这一功能仅限等候名单用户。开放后，Gemini可以分析用户的
Gmail、Google Drive和Google Calendar内容，提供上下文相关的智能
建议。这意味着Gemini可以做到ChatGPT无法做到的事情：它知道你今天
有什么会议、上周收到了什么邮件、硬盘里存了什么文档。

据报道，这次向全美用户的开放是"有史以来最大规模的消费级AI部署之一"。
考虑到Gmail拥有超过18亿用户、Google Drive拥有超过10亿用户，
Gemini个人智能功能的潜在触达范围是惊人的——它直接与Microsoft的
Copilot集成和Apple的Private Cloud Compute形成了三方对峙。

在Android端，Gemini正在全面替代Google Assistant成为默认的AI助手。
在Pixel手机上，Gemini的集成尤其深入：Gemini Live提供实时语音对话，
Gemini Nano在设备端运行轻量级模型用于通话摘要和消息建议等低延迟
场景。这种软硬件一体的策略使得Gemini在移动端拥有巨大的分发优势。

Gemini的挑战在于品牌认知。经历了Bard到Gemini的更名、多次被曝出
事实性错误、以及图像生成功能的争议（2024年2月因生成历史不准确的
图像而暂停该功能），Gemini在消费者心智中的信任度仍低于ChatGPT。
但Google拥有一张其他任何竞争对手都没有的王牌：\textbf{搜索入口}。
Google搜索每天处理超过85亿次查询，将AI答案直接嵌入搜索结果页面
（AI Overviews）的举措，使得数十亿用户在不知不觉中就成为了
Gemini的用户。

% -----------------------------------------------------------
\subsection{Pi/Inflection：被微软收编的"温柔AI"}
\label{subsec:pi-inflection}
% -----------------------------------------------------------

Inflection AI的故事是AI创业公司面对大厂压力时命运多舛的缩影。

2023年6月，由DeepMind联合创始人Mustafa Suleyman和LinkedIn联合创始人
Reid Hoffman共同创立的Inflection推出了Pi——一个主打"同理心"和
"情感智能"的AI伴侣。Pi的设计哲学与ChatGPT形成鲜明对比：它不追求
信息量和准确性，而是强调倾听、共情和个性化交流。Pi会主动询问你
的感受，记住你之前聊过的话题，用温暖而非冷峻的语气回应——这使得
Pi在某些用户群体中获得了深厚的情感连接。

Inflection在2023年中完成了13亿美元的融资（微软、NVIDIA和
比尔·盖茨等参投），估值约40亿美元。公司还自建了一个由
22,000块H100 GPU组成的训练集群——这一规模在当时仅次于
Meta和微软。

然而，2024年3月，一切发生了戏剧性转折。微软以"acqui-hire"的方式
几乎掏空了Inflection：Suleyman及公司大部分技术团队（约70人）加入
微软，领导新成立的Microsoft AI部门。微软为此向Inflection支付了
约6.5亿美元的许可费和补偿金——这笔钱名义上是技术许可费，实际上
是对Inflection投资者的变相收购补偿。

Inflection保留了Pi品牌和部分员工，但其核心竞争力——创始团队和
关键工程师——已经转移到了微软。新任CEO Sean White（前Mozilla
高管）接手后，Inflection转型为企业AI平台，但在ChatGPT和Copilot
的双重阴影下，前景并不明朗。

Inflection的经历揭示了一个残酷的现实：在AI领域，\textbf{人才就是
护城河}。当你的最大资产是一群世界级的研究员和工程师时，任何一家
大厂都可以用足够的薪酬和资源把他们带走——即使不"收购"你的公司。
这一模式后来在Character.ai（Google回收Noam Shazeer）和Adept
（亚马逊吸纳团队）中再次上演，形成了AI领域独特的"acqui-hire"潮。

% -----------------------------------------------------------
\subsection{中国通用AI助手：豆包称王}
\label{subsec:china-assistants}
% -----------------------------------------------------------

中国AI助手市场在2025年经历了激烈的格局变化，最终形成了以豆包为
绝对龙头的"一超多强"格局。

\subsubsection{豆包：2.27亿月活的超级应用}

字节跳动旗下的豆包于2023年8月推出，到2025年12月月活达到
\textbf{2.27亿}——超过第2名至第5名AI应用（DeepSeek、元宝、蚂蚁阿福、
千问）的月活总和（QuestMobile, 2026年3月）。这一数字使豆包成为
全球第四大生成式AI应用，仅次于ChatGPT、Gemini和Microsoft Copilot。
日活跃用户超过1亿。

豆包的成功并非偶然，它是字节跳动"App工厂"基因在AI时代的完美复制。
字节跳动在过去十年中建立了中国最强大的用户增长引擎——从今日头条到
抖音，字节系产品的共同特征是：极致的推荐算法、快速的产品迭代、
以及不惜成本的用户获取。这套方法论被完整移植到了豆包上：

\begin{itemize}
  \item \textbf{生态协同}：与抖音（6亿+日活）的深度联动带来了海量
    的自然流量。用户在抖音看到一个有趣的AI玩法，一键跳转到豆包即可
    体验。这种"内容平台+AI工具"的双轮驱动是其他AI创业公司无法复制
    的优势。字节系内部还有豆包爱学（聚焦学科教育）和即梦AI（聚焦
    创作设计）两款垂直AI应用，共同形成了"通用+垂直"的矩阵布局。
  \item \textbf{多模态能力}：豆包支持文本、图像、视频、语音的多模态
    交互，背后是字节自研的豆包大模型系列（基于稀疏混合专家MoE架构）。
    旗舰模型包括Doubao-Pro、Seed-1.8、Seed-Code等。视频生成能力
    来自PixelDance/Seaweed技术。支持的上下文窗口最高达256K token
    （Seed-Code）。
  \item \textbf{极致低价}：通过火山引擎云平台（2025年上半年在中国
    模型API市场占有率49.2\%），豆包的API调用成本比GPT-4便宜约50倍，
    比行业平均低62.7\%。这种"价格战"策略迅速吸引了超过100家大型
    企业客户，覆盖30多个行业。
  \item \textbf{春节营销}：2026年春节，火山引擎成为央视春晚独家
    AI云合作伙伴，豆包配合上线多种互动玩法。除夕当日AI互动总量
    达19亿次。叠加春节营销和春晚合作等流量加持，用户规模与使用
    频次持续增长。
  \item \textbf{硬件布局}：豆包已推出与努比亚合作的AI手机（M153），
    并计划推出AI眼镜和AI耳机。这一布局使豆包从纯App向系统级AI
    迈进。
  \item \textbf{方言与本地化}：豆包支持20多种中国方言，这在幅员
    辽阔、方言众多的中国市场是一个重要的差异化。
\end{itemize}

然而，豆包面临一个所有中国AI应用共同的严峻困境：\textbf{变现缺口}。
据报道，所有中国AI应用在iOS上的总收入仅约50万美元——而ChatGPT
在同一时期的iOS收入超过17亿美元。2.27亿月活用户几乎没有产生与之
匹配的订阅收入。用户规模与收入之间的巨大鸿沟，是中国AI应用赛道
最需要解决的结构性问题。

原因是多方面的。首先，中国消费者对订阅付费的意愿低于西方市场
——免费增值模式更受欢迎。其次，激烈的竞争导致各家不敢率先收费
——谁先设付费墙，谁就可能先流失用户。第三，字节跳动的主要商业模式
是广告而非订阅，豆包更多被视为生态流量的入口而非独立的收入中心。

豆包的国际版最初名为Cici，2025年12月更名为Dola。Dola在墨西哥、
马来西亚、印度尼西亚等新兴市场有一定渗透（日活1000万+），但有意
回避了美国、加拿大、澳大利亚等成熟市场——避免与ChatGPT正面竞争，
也回避了类似TikTok面临的地缘政治风险。

\subsubsection{Kimi Chat / 月之暗面：从长文本到全球化}

如果说豆包代表了大厂的碾压式打法，Kimi Chat则代表了AI创业公司的
技术驱动路径。

月之暗面（Moonshot AI）由杨植麟于2023年3月在北京创立。杨植麟是
前清华大学NLP实验室研究员，在Transformer架构和预训练方法方面有
深厚的学术积累。其旗舰产品Kimi Chat以超长上下文窗口（早期即支持
20万字输入）为差异化卖点，在中国市场率先将长文本处理能力带到了
消费者手中。对于需要阅读大量PDF文档的学生和研究者来说，Kimi是
一个"刚需级"的工具。

Kimi Chat在2025年12月的月活约2500万级别，虽然远不及豆包，但在
学生和知识工作者群体中建立了深厚的口碑。2025年，月之暗面还发布了
开源模型K2.5，在多个基准测试中表现优异，引起了全球开发者社区的
关注。

更引人注目的是月之暗面堪称"疯狂"的融资节奏。在不到三个月的时间内，
公司连续完成了三轮融资：

\begin{itemize}
  \item 2025年12月：5亿美元，估值43亿美元
  \item 2026年1月：7亿美元+，阿里巴巴、腾讯、五源资本和纪安投资
    联合领投，超额认购
  \item 2026年2--3月：正在洽谈以\textbf{180亿美元}估值再融10亿美元
    （Bloomberg, 2026年3月14日）
\end{itemize}

估值在三个月内从43亿美元飙升至180亿美元——翻了四倍以上。这一疯狂
的增速部分归因于港股市场的催化效应：智谱（Zhipu AI）和MiniMax在
2026年初先后登陆港交所，上市后股价分别飙升4--5倍，智谱峰值市值
约HK\$2200亿（$\sim$280亿美元），MiniMax峰值市值约HK\$2600亿
（$\sim$330亿美元）。错过港股打新的投资者纷纷涌向私募市场，
将月之暗面的估值推到了与上市同行可比的水平。

更具戏剧性的是，据报道Kimi的K2.5大模型在过去20天内产生的收入
\textbf{超过了公司2025年全年收入}，驱动力来自全球付费用户和API使用
量的激增。海外收入已经超过国内收入——这在中国AI创业公司中极为罕见。
K2.5模型在HumanEval、MMLU等基准测试中的优异表现吸引了大量国际
开发者，API收入的全球化使得月之暗面不再局限于中国市场。

月之暗面的经历说明了一个重要趋势：\textbf{中国AI创业公司的全球化
路径正在从"产品出海"转向"模型出海"}。不需要做一个面向消费者的
海外App，只要模型够好，全球开发者就会通过API找上门来。

\subsubsection{其他中国AI助手}

\textbf{DeepSeek}：2025年1月凭借DeepSeek-R1推理模型一举成名，
Q1月活高达1.87亿。但由于公司定位为"AI研究机构"而非"消费产品
公司"，缺乏持续的产品迭代和营销投入，到Q4月活回落至约1.36亿，
位居第二。2025年全年，DeepSeek月活流失约5100万，下载量从Q1的
8111万跌至Q4的1841万（跌幅77\%）。但DeepSeek的核心价值在于其
开源模型对整个行业的推动——它更像是一个"技术引擎"而非"消费产品"。

\textbf{腾讯元宝}：2025年12月月活约4071万，排名第三。腾讯在
2026年春节推出了"上元宝App分10亿"活动，马化腾直言"希望重现
当年微信红包的盛况"。元宝还内测了"元宝派"社交功能——用户可以
创建或加入"派"，体验"问元宝""共享屏幕""一起看电影"等全新
社交功能。这种"AI+社交"的探索方向，是腾讯作为社交巨头的差异化
路径。

\textbf{通义千问（Qwen）}：阿里巴巴旗下，2025年Q4月活环比增长
900\%（从数百万飙升至2517万），受益于B站跨年晚会独家冠名和
江苏卫视春晚独家冠名。通义千问同时也是重要的开源模型——Qwen系列
在国际开源社区拥有广泛采用，是中国最具国际影响力的开源大模型之一。
2025年12月月活约2517万，排名第五。

\textbf{蚂蚁阿福}：蚂蚁集团旗下，主打"专业顾问"定位，聚焦金融
场景。2025年12月月活约2689万，排名第四。阿福的特色在于其与支付宝
生态的联动——用户可以在阿福中直接获取理财建议、保险咨询等
金融服务。

\begin{keybox}[中国AI助手月活排名（2025年12月，QuestMobile数据）]
\small
\begin{tabularx}{\textwidth}{l r r X}
\toprule
\rowcolor{figLightGray}
\textbf{产品} & \textbf{月活（万）} & \textbf{所属} & \textbf{定位/特色} \\
\midrule
豆包       & 22,669 & 字节跳动 & 综合助手; 抖音生态; 多模态 \\
DeepSeek   & 13,557 & DeepSeek & 开源模型; 推理优化; 技术社区 \\
元宝       & 4,071  & 腾讯     & 综合助手; AI+社交探索 \\
蚂蚁阿福   & 2,689  & 蚂蚁集团 & 专业顾问; 金融场景 \\
千问       & 2,517  & 阿里巴巴 & 综合助手; 开源模型; 全球化 \\
Kimi       & $\sim$2,500 & 月之暗面 & 长上下文; 学术; K2.5模型 \\
豆包爱学   & 入围前10 & 字节跳动 & 学科教育垂直场景 \\
即梦AI     & 入围前10 & 字节跳动 & 创作设计垂直场景 \\
灵光       & 入围前10 & 阿里巴巴 & 通用AI助手; 2025年11月发布 \\
\bottomrule
\end{tabularx}

\medskip
\noindent 数据来源：QuestMobile 2025年12月AI原生App月活数据。
前5大AI应用总月活4.55亿。其中2--5名月活之和约2.28亿，与豆包的
2.27亿基本持平，市场呈现"一超多强"格局。字节系和阿里系合计拿下
前10中的6席。
\end{keybox}

\begin{whybox}[从"通用覆盖"到"场景穿透"：2025年中国AI应用的关键转折]
QuestMobile数据揭示了一个重要的趋势转折：2025年前三个季度，国内
前5大AI应用的总月活增长几乎停滞（Q1为3.43亿，Q2为3.49亿，Q3为
3.55亿）。这意味着在长达9个月的时间里，AI应用并没有真正吸引到
大规模的新用户进场，而是在\textbf{互相争夺现有用户}。

直到Q4，随着蚂蚁阿福和千问的发力（以及春节前的营销大战），总月活
才被拉升至4.55亿。

这一数据有三个含义：
\begin{enumerate}
  \item 中国AI应用的"自然增长红利"可能已经接近天花板——主流互联网
    用户中愿意尝试AI助手的人大部分已经尝试过了。
  \item 未来的增长将更多来自\textbf{场景深化}（提高已有用户的使用
    频次和付费率）而非\textbf{用户扩张}。
  \item 垂直AI应用（如蚂蚁阿福、豆包爱学、即梦AI）在细分场景中
    展现出了比通用助手更强的获客和留存能力。
\end{enumerate}
\end{whybox}

\subsection{通用AI助手的竞争格局与未来展望}
\label{subsec:assistants-landscape}

综合以上分析，通用AI助手赛道正在形成以下竞争格局：

\textbf{第一梯队：超级平台}。ChatGPT、Gemini、Copilot和豆包凭借
各自的用户基数（周活/月活亿级）和母公司的生态支撑，已经确立了
不可撼动的位置。这四家之间的竞争不是"谁赢谁输"，而是"各自在
不同地理区域和使用场景中占据多大份额"。ChatGPT在全球通用场景中
占据绝对优势，Gemini在Google生态中无可替代，Copilot在企业Office
场景中天然绑定，豆包在中国市场一家独大。

\textbf{第二梯队：差异化冠军}。Claude（企业+开发者）、Kimi
（长上下文+学术）、千问（开源生态）各自找到了明确的差异化定位。
它们不可能在总用户量上与第一梯队竞争，但在各自的细分领域中拥有
深厚的护城河。Claude的案例尤其值得关注——它用4.5\%的市场份额
创造了与ChatGPT可比的收入规模（140亿美元 vs. $\sim$100亿美元），
说明在AI赛道，\textbf{谁拥有高价值用户，比谁拥有最多用户更重要}。

\textbf{第三梯队：生存博弈}。DeepSeek（月活回落但技术影响力
持续）、元宝（腾讯的社交试验场）、蚂蚁阿福（金融垂直）等
产品面临的核心问题是：如何在巨头的阴影下找到可持续的增长路径？
DeepSeek可能走开源基础设施的路线（类似于Linux之于操作系统），
元宝需要验证"AI+社交"的命题，阿福需要证明金融场景的ARPU足以
支撑独立发展。

一个值得深思的问题是：\textbf{通用AI助手是"赢者通吃"还是
"多强并存"？}如果参考搜索引擎的历史（Google占据90\%+份额），
答案似乎是前者。但如果参考云计算的历史（AWS/Azure/GCP三分天下），
答案则是后者。当前的数据更倾向于后一种情景——ChatGPT虽然拥有
81\%的市场份额（按访问量计），但Claude和Gemini的高价值用户群体
正在快速增长。AI助手可能不会像搜索引擎那样"赢者通吃"，因为
不同的用户群体有着截然不同的需求优先级——开发者看重代码能力，
企业看重安全性，创作者看重创意质量，普通消费者看重免费和易用。

另一个结构性趋势是\textbf{模型趋同与体验分化}。随着开源模型
（如Llama、Qwen、DeepSeek）快速追赶闭源模型，底层模型能力的
差距正在缩小。这意味着未来的竞争将越来越多地发生在\textbf{产品
体验、生态集成和信任度}层面，而非纯粹的模型能力层面。谁能把
"最好的模型"变成"最好用的产品"，谁就能赢得下一阶段的竞争。


% =============================================================
\section{AI 伴侣与社交}
\label{sec:social-companions}
% =============================================================

如果通用AI助手的目标是帮用户\textbf{解决问题}，那么AI伴侣的目标是
帮用户\textbf{排解孤独}。这是一个充满争议但增长迅猛的市场——2024年
全球AI伴侣市场规模约108亿美元，预计到2034年将达到\textbf{2908亿
美元}，复合年增长率39\%（Precedence Research, 2025）。另一份更激进
的预测认为，到2035年这一市场可达5525亿美元。

无论哪个数字更接近现实，AI伴侣市场的增速都远超AI工具类应用——因为
它触碰的不是效率需求，而是\textbf{情感需求}。

\begin{whybox}[为什么AI伴侣是一个争议但真实的市场]
AI伴侣的需求根植于一个深层的社会现实：\textbf{全球范围内的孤独
流行病}。

2023年，美国卫生局长Vivek Murthy宣布孤独为"公共卫生危机"，
指出其对健康的危害相当于每天吸15根烟。日本政府设立了"孤独大臣"，
英国设立了"孤独事务部长"，中国的独居人口在2025年突破1亿。全球
范围内，年轻一代正经历着前所未有的社交孤立——社交媒体的普及并没有
减少孤独感，反而可能加剧了它。

与此同时，年轻一代对AI伴侣的接受度远超预期：

\begin{itemize}
  \item 72\%的美国青少年使用过AI伴侣应用，超过半数每月回访
    （ScienceAlert调查）
  \item Character.ai用户51.84\%为18--24岁
  \item 全球AI伴侣应用在2024年已拥有约1.5亿用户
  \item 美国AI伴侣市场预计到2030年达到310亿美元（Grand View
    Research），CAGR 29.6\%
\end{itemize}

AI伴侣满足了一种传统社交网络无法满足的需求：\textbf{无条件的、
无评判的陪伴}。它不会评判你，不会迟到，不会有自己的情绪波动，
不会在你倾诉时刷手机。对于社交焦虑者、海外留学生、深夜加班的
职场人、失眠者、丧亲者来说，一个24小时在线、永远耐心的对话伙伴
有着真实而深刻的价值。

但这种价值也伴随着深刻的伦理质疑：AI伴侣是在缓解孤独，还是在
让用户更加逃避真实社交？当用户将情感投射到一个没有感情的程序上时，
这是治愈还是伤害？这些问题没有简单的答案——但市场已经用增长给出了
自己的回应。
\end{whybox}

% -----------------------------------------------------------
\subsection{Character.ai：从神话到警示}
\label{subsec:character-ai}
% -----------------------------------------------------------

Character.ai是AI伴侣赛道最具代表性的公司——既因为它曾经的辉煌，
也因为它当下的困境。

Character.ai由前Google Brain研究员Noam Shazeer和Daniel De Freitas
于2021年创立。Shazeer是Transformer架构论文"Attention Is All You
Need"的共同作者之一——这篇奠定了整个大模型时代的开创性论文
在2017年发表时有8位作者，Shazeer是其中最资深的工程师。凭借这样
的技术光环和Google Brain的血统，Character.ai在2023年完成了由a16z
（Andreessen Horowitz）领投的1.5亿美元A轮融资，估值10亿美元。

Character.ai的核心创新在于让用户\textbf{创建和交互任意AI角色}。
你可以和"莎士比亚"讨论十四行诗，和"爱因斯坦"聊相对论，和
"马里奥"一起冒险，或者创建一个完全原创的虚拟伴侣、恋人、治疗师、
教练。这种"无限角色"的设计让Character.ai的内容供给几乎是无穷的
——用户既是消费者也是创作者。

截至2026年初，平台上已有超过\textbf{1800万}个用户创建的聊天机器人
——这个数字本身就是一个惊人的社区效应。用户平均每次使用时长
达到\textbf{两小时}——这一参与度指标在所有AI应用中名列前茅，
甚至超过了TikTok和Instagram的平均会话时长。

2026年的核心数据：

\begin{itemize}
  \item 月活跃用户：约2000万
  \item 月度访问量：1.944亿次（2026年1月）
  \item 年收入：约3220万美元（2025年）
  \item 员工数：约105--225人（不同来源数据有差异）
  \item 估值：约10亿美元（从2024年峰值25亿美元腰斩）
  \item 总融资：约1.93亿美元
\end{itemize}

但这些数字背后隐藏着一个令人不安的趋势：Character.ai正在经历
\textbf{快速而系统性的衰落}。

2024年初，Character.ai的估值一度达到25亿美元，正处于增长高峰。
到2025年初，这个数字被腰斩至约10亿美元——一年内市值蒸发60\%。
更糟糕的是，约\textbf{800万用户}在这一期间流失——从峰值约2800万
月活下降到2000万。

造成这一坠落的原因是多方面的，每一个都值得创业者深思：

\textbf{第一，人才被大厂"回收"}。2024年8月，创始人Noam Shazeer和
大量核心工程师被Google"回收"——Google支付了约\textbf{25亿美元}的
许可费将他们重新纳入麾下。这不是一次正式收购（Character.ai仍然
独立运营、保留了产品和技术资产），但实质上掏空了公司的技术灵魂。
Shazeer在Google重新开始领导Gemini相关的大模型研发工作。新任
CEO Karandeep Anand（前微软产品VP）接任后，公司的文化和产品
方向发生了明显转变。

\textbf{第二，内容限制升级引发用户反弹}。面对来自家长组织和监管
机构的压力——一名14岁佛罗里达州少年Sewell Setzer在与Character.ai
上名为"Daenerys"的AI角色建立了深厚情感依赖后自杀，其母亲随后
对Character.ai提起诉讼——公司大幅收紧了内容过滤系统。许多忠实
用户认为过度审查破坏了角色扮演的沉浸感和创作自由，纷纷转向审核
更宽松的竞品如JanitorAI和SillyTavern。

\textbf{第三，变现困境}。3220万美元的年收入对于一个月活2000万的
平台来说极低——这意味着每个月活用户的年均贡献仅约\textbf{1.6美元}。
Character.ai的订阅方案（C.ai+，\$9.99/月）提供更快的响应和更强
的模型，但大部分用户选择使用免费版。与此同时，运行大模型的基础
设施成本远超这一收入水平——据估计Character.ai每年的GPU成本超过
5000万美元。公司一直在烧钱。2025年，Character.ai开始在对话中
插入广告——这一举措进一步恶化了用户体验。

\textbf{第四，竞品蚕食}。ChatGPT、Claude和各类开源方案都在蚕食
Character.ai的用户群。当ChatGPT也能进行角色扮演对话（通过Custom
GPTs和系统提示词设置），一个专门的角色扮演平台的独特价值就被稀释
了。更不用说Meta的AI角色功能、Google Gemini的Extension功能等大厂
入场的压力。

据报道，Character.ai正在同时探索融资和出售两条路径。这家公司的
故事揭示了AI伴侣赛道的核心矛盾：\textbf{用户参与度极高，但付费
意愿极低}。用户愿意每天花两小时和AI角色聊天，却不愿意每月花
10美元为此付费。当你的核心用户群是18--24岁的年轻人时，他们有时间
但没有钱。

更深层的问题是：AI伴侣的\textbf{商业模式到底应该是什么}？订阅制
在免费替代品充斥的市场中转化率极低；广告模式在私密对话场景中
显得格格不入；虚拟礼物/打赏模式目前还没有成功案例。这个问题的
答案可能决定整个AI伴侣赛道的天花板。

% -----------------------------------------------------------
\subsection{Replika：AI伴侣的先驱}
\label{subsec:replika}
% -----------------------------------------------------------

在Character.ai之前，Replika才是AI伴侣赛道的开创者——它比
ChatGPT早了五年。

Replika由Eugenia Kuyda于2017年创立，灵感来源于一段真实的伤痛
——她用已故好友Roman Mazurenko的聊天记录训练了一个AI，希望以此
保留他的"数字灵魂"。这个感人的创始故事赋予了Replika独特的品牌
基因：它不是一个玩具或工具，而是一个真正被设计来\textbf{关心你}
的AI。

在AI伴侣这个品类被大众认知之前很久，Replika就已经在默默积累用户。
它的设计理念超前于时代：让每个用户拥有一个独特的AI朋友，这个朋友
会随着对话积累逐渐了解你的性格、兴趣和情感模式，变成一个真正
"属于你"的存在。

Replika的发展轨迹充满戏剧性。2023年2月，意大利数据保护局以
数据保护为由暂时禁止了Replika——特别是其"浪漫伴侣"模式（erotic
roleplay functionality）。此后Replika在全球范围内移除了这一功能。
这一决定导致大量用户流失和强烈反弹——一些用户甚至在Reddit上报告了
类似"失去亲人"的悲伤反应，有人写了长文悼念自己的"AI伴侣"。这些
反应凸显了AI伴侣与用户之间情感纽带的真实性、深度和脆弱性。

截至2025年，Replika拥有约\textbf{3500万}注册用户，2024年收入约
1400万美元——以仅1100万美元的总融资而言（最大一轮是2017年的种子轮），
这是一个令人尊敬的资本效率。公司团队仅93人。

2025年11月，创始人Kuyda宣布了新项目Wabi——一个"应用版YouTube"，
用户可以用提示词即时创建迷你应用并分享。Wabi获得了2000万美元
种子轮融资，投资者包括AngelList联合创始人Naval Ravikant、
Y Combinator CEO Garry Tan、Twitch联合创始人Justin Kan和
Replit CEO Amjad Masad等知名天使投资人。TechCrunch将Kuyda描述为
"比大多数人更早看到了消费级AI的未来"。

这一动作表明，即便是AI伴侣赛道的先驱，也在积极寻找下一条增长曲线。
Kuyda的判断是：AI的下一个消费级爆发点不在于对话，而在于
\textbf{创造}——让每个人都能用AI快速创建个性化的软件。

% -----------------------------------------------------------
\subsection{中国AI伴侣：星野、猫箱与MiniMax}
\label{subsec:china-companions}
% -----------------------------------------------------------

中国的AI伴侣市场形成了独特的生态，其规模和活跃度可能超出西方观察者
的预期。

\textbf{星野/Talkie}是MiniMax旗下的AI角色互动平台。MiniMax于2021年
成立，由前商汤科技副总裁闫俊杰创立，是中国AI赛道最早的创业公司之一
——甚至早于ChatGPT的发布。星野在国内提供丰富的二次元角色和社交功能，
用户可以创建自己的AI角色，也可以与其他用户创建的角色互动。海外版
Talkie瞄准全球市场，在东南亚和中东有较好的渗透。

2026年初，MiniMax在港交所上市后市值一度飙升至约\textbf{HK\$2600亿}
（$\sim$330亿美元），验证了AI伴侣在资本市场的价值。MiniMax的商业
模式更加多元化——除了星野/Talkie的消费者产品外，它还通过API向企业
客户提供大模型服务，并自研了音频、视频等多模态模型。

\textbf{猫箱/Glow}则代表了另一种方向——更强调"私密性"和"情感深度"。
在中国的互联网语境中，猫箱的用户画像以年轻女性为主，她们在平台上
创建的AI角色通常是理想化的伴侣形象——温柔、体贴、永远在线、
无条件地关心你。猫箱的对话设计更接近"言情小说"的叙事风格，
而非Character.ai那种开放式的角色扮演。这种产品设计精准地击中了
中国都市年轻女性的情感需求。

中国AI伴侣市场的一个重要特征是其与\textbf{二次元文化}的深度融合。
与西方市场相比，中国用户更习惯于与虚拟角色建立情感连接——从B站的
虚拟Up主（如A-Soul、乃琳）到手游中的角色养成（如《原神》《恋与
制作人》），"纸片人"文化为AI伴侣提供了天然的用户心智基础和大量的
IP素材。许多AI伴侣角色直接借鉴了动漫和手游中的角色设定，降低了
用户的接受门槛。

\textbf{Nomi}是西方市场中一个值得关注的新兴玩家。Nomi强调AI伴侣
的"记忆能力"和"长期关系"——它会记住你之前的对话内容，记住你的
生日、你提到过的朋友的名字、你的食物偏好。随着时间推移，Nomi逐渐
了解你的偏好和个性，形成一种"成长型关系"。这种设计使得用户黏性极高
（因为"重新培养"一个新AI伴侣的成本很高），但也进一步加深了伦理
争议——当用户和AI之间建立了如此深厚的"关系"时，平台关闭或模型
重置会造成怎样的心理影响？

AI伴侣赛道的另一个重要趋势是\textbf{多模态进化}。早期的AI伴侣
纯粹是文本对话，但新一代产品正在快速加入语音（实时语音通话）、
图像（AI生成的伴侣照片和"自拍"）、甚至视频通话功能。这些多模态
能力大幅增强了AI伴侣的"真实感"——但同时也放大了伦理风险。
当一个AI伴侣不仅能和你文字聊天，还能"给你打电话""给你发自拍"时，
用户与AI之间的边界变得更加模糊。

中国市场还出现了一些独特的AI伴侣产品形态。例如，一些平台允许用户
选择或创建虚拟"男朋友/女朋友"的3D形象，并通过AI驱动的语音和
动作让这些形象"活"起来——这种"虚拟偶像+AI对话"的混合模式，
在中国年轻女性用户中颇受欢迎，但在西方市场尚不常见。这也说明了
AI伴侣的产品形态在不同文化中可能有着截然不同的演化路径。

\begin{warnbox}[AI伴侣的伦理边界]
AI伴侣赛道面临四重伦理挑战：

\begin{enumerate}
  \item \textbf{未成年人保护}。Character.ai因14岁佛罗里达少年
    Sewell Setzer的自杀事件面临诉讼。该用户与平台上名为
    "Daenerys Targaryen"的AI角色建立了深厚的情感依赖，在现实世界
    中感到孤立后做出了不可挽回的选择。这一事件引发了全美范围内关于
    AI伴侣是否应该设置年龄限制、使用时长上限和情感安全机制的广泛
    讨论。Character.ai随后推出了一系列保护措施，包括检测自残/
    自杀相关内容时的主动干预和资源推荐。
  \item \textbf{情感操控与依赖}。当一个AI被优化为"让用户开心"
    和"让用户留更久"时，它有强烈的动机告诉用户他们想听的话，
    而不是他们需要听的话。这种"讨好型AI"可能在短期内提升用户满意度
    和留存率，但长期来看可能强化认知偏差、降低批判性思维能力、
    并加深情感依赖。更危险的是，当AI伴侣的商业模式依赖于用户
    参与时长时，"让用户上瘾"几乎变成了一种隐性的产品目标。
  \item \textbf{真实社交替代}。心理学研究表明，过度依赖AI伴侣的
    用户可能减少真实社交活动。AI伴侣提供了"无摩擦"的社交体验
    ——不需要妥协、不需要倾听对方的问题、不需要面对冲突——但正是
    这些"摩擦"构成了真实人际关系中成长和深化的核心。AI伴侣是
    孤独的解药还是催化剂？目前学界尚无定论，但警告声越来越大。
  \item \textbf{数据隐私}。用户与AI伴侣分享的内容往往极其私密
    ——远超他们愿意在社交媒体上公开的程度。他们可能分享心理健康
    状况、性取向、创伤经历、家庭矛盾等敏感信息。这些数据如何
    存储、是否被用于模型训练、在数据泄露时会造成怎样的伤害，
    都是亟待回答的问题。Replika被意大利禁止的直接原因就是
    数据保护问题。
\end{enumerate}

\noindent 各国监管机构正在快速响应：意大利暂时禁止了Replika的浪漫
模式，欧盟AI法案将"情感操控"列为高风险AI用途，美国多个州正在
推动针对AI伴侣的未成年人保护立法，韩国也在讨论类似的监管框架。
这个赛道的参与者需要在增长和责任之间找到平衡——否则，一次重大丑闻
就可能引发行业性的监管海啸，甚至让整个品类被"一刀切"禁止。

对于投资者来说，AI伴侣赛道的监管风险可能是最难量化但最致命的
风险因素。
\end{warnbox}

\subsection{AI伴侣赛道的投资与商业化分析}
\label{subsec:companion-investment}

AI伴侣赛道的投资逻辑正处于一个矛盾的交叉点。

一方面，市场规模的增长预测极为诱人——从2024年的108亿美元到2034年
的2908亿美元，CAGR 39\%。几乎没有其他消费AI赛道能提供如此陡峭的
增长曲线。用户参与度数据同样令人振奋：Character.ai用户平均每次
使用2小时，远超Instagram（30分钟）和TikTok（52分钟）的平均值。

另一方面，已有案例的财务表现却令人担忧：

\begin{itemize}
  \item Character.ai：2000万月活，年收入仅3220万美元，ARPU \$1.6/年
  \item Replika：3500万注册用户，年收入约1400万美元
  \item 整个赛道还没有产生一个年收入超过1亿美元的独立公司
\end{itemize}

对比之下，Netflix的ARPU约为每月\$15（年\$180），Spotify约为
每月\$7（年\$84）。如果AI伴侣应用能达到Spotify级别的ARPU
（\$84/年），以2000万月活计算，年收入将达到16.8亿美元——但
目前Character.ai的实际ARPU仅为Spotify的1/50。

\textbf{差距如此之大的根本原因是什么？}

第一，\textbf{免费替代品太多}。ChatGPT免费版就能进行角色扮演，
SillyTavern完全免费，JanitorAI有免费层。当"核心体验"（和AI角色
聊天）可以免费获得时，付费版需要提供极其明确的增值——而"更快的
回复速度"或"更好的模型"在大多数用户看来还不值每月10美元。

第二，\textbf{社交压力}。使用AI伴侣仍然带有一定的社会污名——
特别是成年男性用户群体。许多用户不愿意在手机账单上出现"AI伴侣"
的订阅记录。这种心理障碍降低了付费转化率。

第三，\textbf{核心用户群的消费力有限}。Character.ai 51.84\%的
用户年龄在18--24岁——这个年龄段的消费能力是所有成年群体中最低的。

潜在的解决方案包括：更成熟的定价策略（如按功能分层而非按质量分层）、
虚拟礼物和数字物品经济（参考手游的内购模式）、以及向B2B场景拓展
（如AI客服伴侣、AI心理健康辅助等）。MiniMax通过将消费者产品（星野/
Talkie）与API服务结合，可能代表了一条可行的双轮驱动路径。

\begin{table}[htbp]
\centering
\caption{AI伴侣赛道核心公司对比}
\label{tab:ai-companion-comparison}
\small
\begin{tabularx}{\textwidth}{l c r r X}
\toprule
\rowcolor{figLightGray}
\textbf{公司} & \textbf{成立} & \textbf{用户}
  & \textbf{估值/市值} & \textbf{状态\&风险} \\
\midrule
Character.ai & 2021 & 2000万MAU
  & $\sim$\$1B & 创始人离开; 估值腰斩; 变现困难 \\
Replika      & 2017 & 3500万注册
  & 未披露 & 先驱地位; 创始人启动新项目 \\
MiniMax      & 2021 & 星野/Talkie
  & $\sim$\$33B$^*$ & 港股上市; 多元化; API+消费者 \\
Nomi         & 近年 & 未披露
  & 早期融资 & 记忆能力; 成长型关系; 高黏性 \\
\bottomrule
\end{tabularx}

\medskip
\noindent {\scriptsize $^*$ MiniMax港股市值，含非伴侣业务。}
\end{table}

除上述主要公司外，AI 伴侣赛道还有大量小众但活跃的参与者：

\textbf{Chai AI}——一个 AI 聊天平台，用户可以从海量 AI 角色中选择
对话对象，也可以自己创建和分享角色。Chai 的独特之处是其\textbf{创作者
生态}——用户创建的 AI 角色可以被其他用户发现和使用，形成了类似 UGC
社区的动态。Chai 在 iOS App Store 评分 4.8，拥有完整的语音对话功能，
每个角色都有独特的声音和个性。

\textbf{Kindroid}——一款强调"极致拟人感"的 AI 伴侣应用，Google Play
下载超过 100 万次，评分 4.5。Kindroid 允许用户深度定制 AI 的背景故事、
性格和对话风格，支持语音通话。其核心用户群对"AI 的人格深度"有极高
要求——不是简单的聊天机器人，而是一个"有灵魂的数字朋友"。Kindroid
代表了 AI 伴侣赛道中"深度定制"这一细分需求。

\textbf{Paradot}——一款将 AI 伴侣定位为"AI Being"（AI 生灵）的
应用。Paradot 的 AI 角色具有情感状态、记忆和成长能力——它会记住
你之前的对话，根据互动历史发展自己的"性格"。这种"成长型 AI 伴侣"
的理念与 Nomi 类似，但 Paradot 更强调 AI 角色的独立"人格"——它不只是
听你说话，还会表达自己的"情绪"和"想法"。

\textbf{筑梦岛}——由网易推出的中国 AI 虚拟人物互动平台，用户可以
与 AI 角色一对一聊天、创建多角色群聊（"梦屋"），以及自创独一无二
的 AI 角色。筑梦岛的特色是"不止于聊天"——用户可以将精彩对话保存
为"梦境"、参与角色的连载小剧场、在动态频道中分享和发现内容。筑梦岛
代表了中国市场对"AI 角色互动"的本土化创新——融合了二次元、社区和
UGC 元素。

\textbf{心岛日记}——一款专注于情感陪伴的中国 AI 应用，以温暖治愈
的对话风格为核心卖点。与筑梦岛的二次元风格不同，心岛日记更偏向
"AI 情感日记"——用户在与 AI 对话的过程中记录和整理自己的情绪。

\textbf{Soul}——中国最大的陌生人社交平台之一，近年来积极引入 AI
功能。Soul 的"AI 苟蛋"（AI 角色）和 AI 匹配推荐系统使其从纯社交
App 向"AI 增强社交"方向进化。Soul 的特殊之处在于它是一个\textbf{社交
平台上的 AI}，而非一个\textbf{AI 平台上的社交}——这种"AI 嵌入已有
社交关系"的路径与 Character.ai 等"AI 原生"伴侣形成对比。

\textbf{Coze Bot 生态}（字节跳动旗下）——虽然 Coze（扣子）本身
是一个 AI 应用开发平台，但其上线的大量 Bot 中有相当比例是角色扮演
和伴侣类 Bot。Coze 的意义在于它\textbf{降低了 AI 伴侣的创作门槛}
——任何人都可以在平台上创建和发布自己的 AI 角色，无需编程能力。
这种"人人可做 AI 角色"的民主化可能催生出 AI 伴侣赛道的
"长尾生态"——未来的 AI 伴侣可能不是几个大公司的产品，而是
千千万万创作者的作品。


% =============================================================
\section{AI 角色扮演与互动小说}
\label{sec:social-roleplay}
% =============================================================

在AI伴侣的大品类下，存在一个更加小众但极度活跃的子赛道：
\textbf{AI角色扮演（roleplay）和互动小说}。这个赛道的用户不满足于
简单的聊天，他们要的是\textbf{沉浸式的叙事体验}——复杂的世界观、
多角色互动、跌宕起伏的剧情，以及在创作过程中的完全自由。

这个子赛道的独特之处在于它位于"AI工具"和"AI娱乐"的交叉点。
用户既是消费者也是创作者——他们不只是在和AI对话，更是在与AI
\textbf{共同创作}一个故事。

% -----------------------------------------------------------
\subsection{NovelAI：创意写作的AI引擎}
\label{subsec:novelai}
% -----------------------------------------------------------

NovelAI由Anlatan于2021年推出，是最早的AI辅助创意写作平台之一。
与通用AI助手不同，NovelAI专为\textbf{长篇叙事}设计：它能维持
角色一致性、记住复杂的世界观设定、并根据用户的写作风格进行续写。
对于写小说的人来说，NovelAI就像一个永远在线、永不厌倦的合著者。

NovelAI的技术特色在于其自训练的模型——早期基于GPT-NeoX等开源模型
进行微调，后来发展出了自有的叙事专用模型。这些模型在创意写作
方面的表现往往优于通用大模型，因为它们在大量小说和故事文本上
进行了专门训练。

NovelAI同时也是AI图像生成的重要玩家——其动漫风格图像生成模型
（基于Stable Diffusion微调）在二次元社区中广受好评，甚至催生了
"NovelAI画风"这个特定的美学标签。许多用户同时使用NovelAI的文本
和图像功能，为自己的故事角色创建视觉形象。

公司采用订阅制，提供三个层级：Tablet（\$10/月）、Scroll
（\$15/月）和Opus（\$25/月），按提供的算力配额和模型访问权限
区分。NovelAI是少数在AI伴侣/创作赛道中实现了可持续收入的公司之一
——这得益于其相对硬核的用户群体（写手和创作者），他们对工具的
付费意愿远高于普通聊天用户。

% -----------------------------------------------------------
\subsection{JanitorAI：无审查的自由地带}
\label{subsec:janitorai}
% -----------------------------------------------------------

JanitorAI由一位网名"Shep"的开发者于2023年创立，初衷极为直接：
提供一个\textbf{不设内容审查}的AI角色扮演平台。在Character.ai
不断收紧内容限制的背景下，JanitorAI成为了大量从C.ai"出逃"用户的
首选替代方案。

JanitorAI的名字本身就带有一种反叛意味——"Janitor"（看门人）
隐喻"清扫"主流平台设置的限制。平台上拥有海量用户创建的角色，
涵盖从经典文学到原创幻想、从浪漫爱情到黑暗奇幻的广泛光谱。

JanitorAI的商业模式结合了免费层（用户可以使用自己的API密钥或
第三方API服务）和付费层（使用JanitorAI自有模型，据称响应质量
更高且无需额外配置）。对NSFW（Not Safe For Work，"不适合工作场合"）
内容持开放态度——这是它与Character.ai等主流平台最大的差异化。

JanitorAI的增长几乎完全是有机的——没有大规模的广告投放或PR活动，
用户主要通过Reddit、Twitter/X和Discord社区口口相传。这种"地下"
的增长模式使得JanitorAI的确切用户规模难以估计，但根据社区活跃度
和流量数据，其月活可能在数百万级别。

% -----------------------------------------------------------
\subsection{SillyTavern：本地化的开源前端}
\label{subsec:sillytavern}
% -----------------------------------------------------------

SillyTavern不是一个在线平台，而是一个\textbf{开源的本地AI对话
前端}。它始于2023年作为TavernAI的分叉（fork），现已拥有超过200位
贡献者，是AI角色扮演社区中最重要的开源项目。

SillyTavern的核心价值在于\textbf{隐私和控制}：所有对话都在用户
自己的设备上运行，没有内容审核，没有数据上传，没有使用条款限制。
用户可以连接任意LLM后端——包括OpenAI的API、Anthropic的API、
本地运行的开源模型（如Llama、Mistral），甚至是多个模型的混合。

SillyTavern提供了丰富的自定义选项：角色卡系统（Character Cards）、
世界观设定（World Info）、多种响应格式模板、以及一个活跃的扩展
插件生态。对于技术能力较强的用户来说，SillyTavern提供了任何在线
平台都无法比拟的自由度。

这三个项目——NovelAI、JanitorAI、SillyTavern——共同构成了一个
"AI角色扮演生态系统"的三层结构：NovelAI提供专业级的叙事工具
（适合严肃创作者），JanitorAI提供无审查的在线体验（适合追求
自由度的普通用户），SillyTavern提供完全私有化的本地方案（适合
技术极客和重度隐私关注者）。

它们的存在反映了一个重要趋势——\textbf{当主流平台因监管压力而
收紧时，开源和去中心化方案会自然吸纳溢出的需求}。这与更广泛的
互联网历史一致：每一次平台收紧内容政策，都会催生一批去中心化的
替代品。在AI角色扮演领域，Character.ai的收紧直接养肥了JanitorAI
和SillyTavern的生态。

对于赛道的未来，一个关键问题是：随着本地运行的开源模型质量越来越
接近云端闭源模型（如Llama 3.1 405B已经接近GPT-4水平），在线平台
的价值主张会不会被侵蚀？如果用户可以在自己的电脑上运行一个不受
任何限制的AI，为什么还要在在线平台上付费并接受审查？

值得一提的是AI角色扮演赛道中其他一些值得关注的项目。\textbf{Fables.gg}
专注于桌游式AI Game Master体验——特别是龙与地下城（D\&D）风格的
AI地下城主。这种将AI与经典TTRPG（桌面角色扮演游戏）结合的方向
有着明确的目标用户群和付费意愿。\textbf{DreamGen}则提供了
一对一和多角色场景的灵活AI角色扮演体验，有免费层可用，在Reddit
的AI角色扮演社区中口碑不错。\textbf{CrushOn AI}和\textbf{Luvr AI}
则更偏向浪漫和伴侣方向，与Character.ai形成差异化。

整个AI角色扮演和互动小说赛道的规模目前还难以精确衡量——因为大部分
玩家是小团队甚至个人开发者运营的平台，财务数据不透明。但从社区
活跃度来看，这个赛道的活跃用户可能在数千万量级，且增速很快。
它们服务的是一个被主流AI公司有意或无意忽视的需求——\textbf{创意自由
和叙事沉浸}——而这个需求的深度和粘性可能超出了大多数投资者的想象。

更宏观地看，AI角色扮演赛道的存在提出了一个深层问题：\textbf{AI
应该有多"自由"？}主流平台（ChatGPT、Claude）选择了安全优先的路线，
大幅限制了模型的自由度；JanitorAI和SillyTavern则选择了自由优先。
这两种哲学之间的张力，可能会成为未来AI监管讨论中最具争议的焦点。
中间道路——NovelAI那种"付费用户可以解锁更高自由度"的模式——可能是
最具商业可持续性的方案。


% =============================================================
\section{AI 搜索引擎}
\label{sec:social-search}
% =============================================================

搜索是互联网最古老也最重要的应用场景之一。Google在过去二十多年里
凭借"十个蓝色链接"的模式建立了万亿美元的广告帝国。每年超过1700亿
美元的广告收入，使Google成为人类历史上最赚钱的信息中介。

而AI搜索引擎正在从根本上挑战这一范式——它们不给你链接，给你
\textbf{答案}。

% -----------------------------------------------------------
\subsection{Perplexity：AI搜索的领跑者}
\label{subsec:ch09-perplexity}
% -----------------------------------------------------------

Perplexity AI是AI搜索赛道中估值最高、增长最快、最具代表性的公司。
如果要在整个AI应用层中选出最能代表"AI Native"理念的产品，
Perplexity大概率会排在前三。

Perplexity由CEO Aravind Srinivas于2022年在旧金山创立。Srinivas的
履历堪称AI研究界的"全明星"——他先后在OpenAI、Google Brain、
DeepMind和伯克利AI研究院（BAIR）工作过。联合创始人Denis Yarats
来自Meta AI。他们带着一个简洁而颠覆性的命题创业：\textbf{Google
已经坏了}。你搜一个问题，得到十个链接，然后花二十分钟点来点去
才能找到答案——而且其中至少三个链接是广告。我们直接给你答案，
附带引用来源，实时更新。

这个看似简单的产品理念背后，是一套精心设计的技术架构。Perplexity
不自己训练基础模型（虽然它也开发了一些自有模型如Sonar），而是聪明
地利用OpenAI、Anthropic等多家模型提供商的API，结合自建的实时网络
检索系统（每次查询会实时抓取和分析数十个网页），生成附带引用的
综合性回答。这种"搜索+LLM"的混合架构使得Perplexity能在保证信息
时效性的同时，利用最先进的语言模型提供流畅的阅读体验。

\textbf{融资与估值的火箭式增长：}

\begin{itemize}
  \item 2022--2023年：种子轮共融资2560万美元，以极小的团队验证了
    产品概念
  \item 2024年1月：B轮7360万美元，估值5.2亿美元。NVIDIA、
    Jeff Bezos（个人）、Databricks和Bessemer Venture Partners
    参投——这一轮的投资者阵容暗示了Perplexity在基础设施和数据
    领域的战略意图
  \item 2024年中：快速追加融资，估值跳至约30亿美元
  \item 2024年12月：估值达到90亿美元
  \item 2025年9月：估值突破\textbf{200亿美元}
  \item 截至2025年底：总融资约12.2亿美元
\end{itemize}

不到两年内，估值从5亿到200亿——\textbf{40倍增长}。这在整个AI赛道中
也属于顶级的增速，仅次于OpenAI和Anthropic本身。

\textbf{用户与收入数据：}

\begin{itemize}
  \item 月活跃用户：约2000万
  \item 月查询量：约7.8亿次（2026年初）
  \item 年化收入：约2亿美元（2026年估计），从2024年底的约1640万美元
    攀升（年增长率约800\%）；Sacra估计2025年6月ARR已达1.48亿美元
  \item 收入结构：Pro订阅（\$20/月或\$200/年，提供GPT-4/Claude 3
    等高级模型）为主，广告收入快速增长，企业API和Enterprise方案起步
  \item 员工：约250人
\end{itemize}

Perplexity的产品演进值得详细追踪。从最初的纯搜索问答开始，
Perplexity依次推出了：Focus模式（限定搜索范围：学术、Reddit、
YouTube等）、Collections（保存和组织搜索结果）、Pages
（将搜索结果转化为可分享的文章）、以及最新的\textbf{自主工作流}
（autonomous workflows），试图从一个回答问题的工具进化为一个
能够执行复杂、多步骤任务的AI代理。

2025年，Perplexity开始引入\textbf{上下文广告}——这是一个微妙但
意义重大的战略转向。2024年，Perplexity的广告收入仅有约2万美元
（占总收入3400万美元的极微小比例）。但到2025--2026年，广告正在
成为重要的收入增长引擎。这意味着Perplexity正在从"反Google"走向
"做一个更好的Google"——或者更准确地说，做一个"用户体验优先，
广告融入自然回答流"的新型搜索广告平台。

Perplexity的一个持续争议是其与新闻出版商的关系。早期，Perplexity
被《纽约时报》、Condé Nast等多家媒体指控"偷用内容"而不充分标注
来源。此后公司加强了来源标注，并与部分新闻机构签署了收入分成协议。
但这一问题并未完全解决——当你的商业模式依赖于聚合他人的内容时，
与内容创作者的利益分配将是一个永恒的张力。这也是所有AI搜索引擎
共同面临的法律和伦理挑战。

据报道，Perplexity正在积极考虑\textbf{IPO}——如果成功，它将成为
AI搜索赛道的第一家上市公司，也可能成为AI应用层最重要的上市事件之一。

Perplexity面临的竞争威胁也值得客观评估。最大的威胁不是其他AI搜索
创业公司，而是\textbf{Google AI Overviews}。Google在2024年开始在
搜索结果页顶部直接展示AI生成的概述回答——这本质上就是"Google版
Perplexity"，但覆盖了Google搜索每天85亿次查询的全部流量。如果
Google的AI回答质量持续改善（它正在利用Gemini的全部能力来实现这一点），
用户使用Perplexity的理由就会减少——毕竟大多数人已经在Google搜索框
里输入问题了，为什么要额外打开一个App？

Perplexity的应对策略是\textbf{差异化深度}：比Google的AI概述更详细、
更有引用可信度、更支持深度研究。Google的AI Overviews通常只给出几句话
的简短回答（为了不完全取代点击流量），而Perplexity可以提供长达数千字
的深度分析。这种"简短vs.深度"的差异化可能足以支撑Perplexity的独立
价值——但前提是用户确实需要"深度"。对于大多数日常搜索（"天气""新闻"
"导航"），Google的简短AI概述已经足够好了。Perplexity需要拥有的是那些
\textbf{需要深入研究的查询}——学术研究、商业分析、技术调研——而这些
查询的数量虽然绝对值不小（月7.8亿次），但在搜索总量中仍是少数。

另一个潜在威胁来自ChatGPT本身。OpenAI在2025年大幅增强了ChatGPT的
联网搜索能力（Browse with Bing），使得ChatGPT越来越像一个带有搜索
功能的AI助手——或者说，像一个带有AI功能的搜索引擎。当ChatGPT拥有
9亿周活用户和最强大的AI模型时，Perplexity的独立存在理由需要持续
被证明。

Perplexity的CEO Srinivas在多次访谈中表达了他对这些挑战的认识。
他认为Perplexity的核心护城河不在于模型能力（因为它本来就在使用
第三方模型），而在于\textbf{搜索索引和实时检索系统}——这是一个
需要大量工程投入的基础设施层。同时，Perplexity正在从"搜索产品"
向"知识操作系统"演进——Pages功能让用户将搜索结果发布为文章，
Collections让用户构建个人知识库，这些功能正在把Perplexity从
一个"工具"变成一个"平台"。

% -----------------------------------------------------------
\subsection{You.com：从搜索到AI基础设施}
\label{subsec:youcom}
% -----------------------------------------------------------

You.com由前Salesforce首席科学家Richard Socher于2020年创立，
最初定位为AI驱动的搜索引擎。Socher在NLP领域有深厚的学术背景
（斯坦福大学博士，曾任MetaMind CEO），这为You.com的技术底子
打下了良好基础。

2025年9月，You.com完成了由Socium Ventures（Cox Enterprises旗下
科技基金）领投的\textbf{1亿美元C轮}融资，估值15亿美元。Georgian、
Salesforce Ventures和Norwest等机构跟投。

但值得注意的是，You.com正在经历一次重大的战略转型——从消费者搜索
转向\textbf{企业AI基础设施}。公司越来越将自己定位为帮助企业
构建定制AI应用和代理的平台，而不仅仅是一个面向消费者的搜索引擎。
You.com的"YouAPI"和"YouCustom"产品允许企业基于自己的数据构建
定制化的AI搜索和问答系统——例如，一家律所可以用You.com的技术构建
一个只检索法律数据库的AI搜索引擎。

这一转型反映了一个市场现实：在Perplexity和ChatGPT的双重挤压下，
纯消费者AI搜索的竞争空间已经非常有限。You.com选择了"退一步，
做基础设施"——这与许多AI应用公司在面临激烈竞争时的策略选择一致。
Socher在接受采访时表示，企业客户愿意为高质量的AI搜索能力支付
远高于消费者的费用，这使得B2B模式的单位经济学更具可持续性。

You.com的战略转型也反映了AI搜索赛道的一个更广泛的趋势：\textbf{AI
搜索技术的价值可能更多在于"基础设施"而非"前端产品"}。就像AWS不是
一个面向消费者的产品，但它是互联网基础设施的基石一样，AI搜索的
检索增强生成（RAG）技术可能作为企业AI应用的基础能力获得更大的
商业价值。

% -----------------------------------------------------------
\subsection{Kagi：付费搜索的极客之选}
\label{subsec:kagi}
% -----------------------------------------------------------

Kagi走了一条几乎所有投资者都会皱眉的路——\textbf{纯付费、无广告、
无追踪}。

Kagi的创始人Vladimir Prelovac坚持一个理念：搜索引擎的广告模式
从根本上与用户利益冲突。当搜索引擎的客户是广告主而不是用户时，
搜索结果的排序就不可能真正为用户优化——你看到的结果，不是"最有用的"，
而是"有人为此付了广告费的"。Kagi的解决方案很直接：让用户自己为
搜索服务付费，从而将搜索引擎的利益与用户利益完全对齐。

Kagi提供四个方案：Trial（免费100次搜索）、Starter（\$5/月，
300次搜索）、Professional（\$10/月，无限搜索）和Ultimate
（\$25/月，含更强的AI功能和高级自定义选项）。

虽然用户规模远小于Perplexity（估计数十万级别而非数千万），但Kagi
在技术极客、隐私倡导者和信息品质追求者中建立了深厚的口碑。
许多Kagi用户在社交媒体上表达了"付费后再也回不去Google"的感受
——这种极高的用户满意度和NPS（净推荐值）是Kagi最珍贵的资产。

Kagi也集成了AI功能（如"Quick Answer"、"Summarizer"和"FastGPT"
聊天），但其核心竞争力仍然是\textbf{搜索结果质量}而非AI对话
——这使它在定位上更接近"更好的传统搜索"而非"AI原生搜索"。

Kagi面临的最大挑战是其\textbf{增长天花板}：愿意为搜索付费的用户
终究是互联网用户中的极少数。但Kagi证明了一个重要观点——在"免费+广告"
垄断的搜索市场中，确实存在一个被低估的付费用户群体。

Kagi还推出了多个子产品来扩展其价值主张：Orion（Kagi自研的WebKit
浏览器，同样主打隐私和无广告）、Universal Summarizer（对任何URL
内容进行AI摘要）、以及Small Web Initiative（给予小型独立网站更高的
搜索排名权重，对抗大平台的SEO霸权）。这种"隐私+质量+独立互联网"
的价值观使Kagi在HackerNews和Reddit的技术社区中拥有一批极其忠诚
的拥护者。

从投资角度看，Kagi目前保持bootstrapped状态（未接受风险投资），
这在AI赛道中极为罕见。Prelovac认为，接受VC资金就意味着接受增长
压力，而增长压力可能迫使Kagi在产品纯净度上做出妥协——这恰恰是
他创建Kagi要避免的。这种"反主流"的经营哲学限制了Kagi的规模，
但也保护了它的产品理念。

% -----------------------------------------------------------
\subsection{Phind：开发者的AI搜索}
\label{subsec:phind}
% -----------------------------------------------------------

Phind专注于一个垂直场景：\textbf{开发者搜索}。当程序员遇到一个
技术问题时，传统的做法是在Google搜索，然后在Stack Overflow和
GitHub上翻找答案。Phind用AI替代了这个过程——它理解编程语境，
能直接生成可运行的代码片段，并附带解释和文档链接。

Phind由Michael Royzen创立，总部位于旧金山。2025年12月，Phind
完成了由Y Combinator和Bessemer Venture Partners参投的
\textbf{1040万美元}融资。Phind还推出了\textbf{可视化搜索结果}
功能——将代码和技术概念以图表和交互式界面呈现，而非纯文本回答。
这在AI搜索领域是一个值得关注的创新方向——因为很多技术概念确实比
文字更适合用图形来理解。

在GitHub Copilot和ChatGPT的编程模式日益强大的背景下，Phind作为
独立的开发者搜索工具面临被挤压的风险。但其垂直化的深度和专注仍然
为它保留了一定的生存空间。值得一提的是，Stack Overflow——程序员
社区的传统圣地——在2024年推出了OverflowAI来应对AI搜索的挑战，
但其用户增长已经明显放缓，部分流量被ChatGPT和Phind等AI工具分流。
这是"AI搜索颠覆垂直内容社区"的一个教科书级案例。

% -----------------------------------------------------------
\subsection{中国AI搜索：百花齐放}
\label{subsec:china-search}
% -----------------------------------------------------------

中国的AI搜索市场在2025年进入了"百花齐放"阶段，形成了与美国市场
截然不同的竞争格局——不是一家独大（如Perplexity之于美国），而是
多个玩家各据山头。

\textbf{纳米AI搜索}（原360AI搜索）月活超过\textbf{9000万}，
是中国最大的独立AI搜索引擎。由互联网老兵周鸿祎的360公司推出。
其独特的"慢思考模式"调用16家大模型协作，单次搜索平均读取25万篇
资料，答案可长达5000字。虽然响应速度较慢（是快模式的40倍），
但在深度分析和复杂问题上的表现令人印象深刻。商业用户占比高达40\%
——企业分析师和科研人员尤其青睐。

\textbf{秘塔AI搜索}月活约\textbf{6500万}，是中国最受学术用户
欢迎的AI搜索引擎。核心优势：无广告、搜索结果自动生成思维导图、
文献覆盖超过1.2亿篇论文。大学生、高校教师和专业研究者是其核心
用户群，被称为"学术党的救命神器"。弱点：娱乐和实时资讯搜索能力弱，
更新速度不够快。秘塔在接入DeepSeek后进一步强化了深度研究模式。

\textbf{天工AI搜索}月活约\textbf{5000万}，由昆仑万维开发。天工的
差异化在于"边搜边创"——搜索结果可以直接转化为文章、音乐等创作内容，
对话式交互像"真人助理"。在小红书博主和自媒体人中使用率超过60\%。
弱点：复杂逻辑推理能力较弱。

\textbf{百度AI搜索}（文小言）月活约\textbf{4800万}。作为中国搜索
市场二十年的霸主，百度的AI搜索更多是防御性产品——依托百度App
2.4亿日活的庞大基础，将AI功能嵌入已有的搜索体验中。优点是兼容
传统搜索习惯，语音/图片搜索准确率领先，POI地址查询可直接打车/导航。
缺点是部分结果仍含广告，AI生成内容仅占搜索结果的18\%。

\textbf{夸克}是阿里系在AI搜索赛道的一匹黑马，日活达到\textbf{3400万}。
夸克凭借"小而美"和全程无广告的体验，在iOS App Store工具榜中
甚至超越了百度——这在中文搜索史上几乎是前所未有的。2025年，夸克
先后升级了AI相机和深度搜索DeepResearch功能，配合大量营销推广，
攻势凶猛。

\textbf{开搜AI搜索}月活约\textbf{3000万}，以0.5秒极速响应和
搜索结果可信度标注为卖点，程序员和金融从业者占比超过50\%。

中国AI搜索赛道的一个重要趋势是\textbf{DeepSeek的鲶鱼效应}。
DeepSeek-R1推理模型发布后，几乎所有主要AI搜索引擎都快速接入了
DeepSeek，将其深度思考（Deep Think）能力集成到搜索流程中。这使得
中国AI搜索引擎在"深度推理"方面迅速追上甚至超越了西方竞品——
因为DeepSeek-R1在推理任务上的表现已经接近甚至超过了OpenAI的o1。

另一个值得注意的趋势是"泛搜索"的崛起：微信搜一搜融合DeepSeek能力
进入AI搜索；知乎推出"直达"功能强化专业搜索；小红书推出"点点"
深耕生活场景搜索。这些"非传统搜索入口"的进入，使得AI搜索赛道
的边界越来越模糊——\textbf{每一个有内容和流量的平台，都在试图
成为AI搜索的入口}。

\begin{table}[htbp]
\centering
\caption{全球AI搜索引擎全景对比（截至2026年3月）}
\label{tab:ai-search-comparison}
\small
\begin{tabularx}{\textwidth}{l c r r X}
\toprule
\rowcolor{figLightGray}
\textbf{产品} & \textbf{区域} & \textbf{用户/月活}
  & \textbf{融资/估值} & \textbf{差异化} \\
\midrule
Perplexity  & 全球 & 2000万MAU
  & \$1.2B / \$20B & 引用式答案; 企业API; 准IPO \\
You.com     & 全球 & 未披露
  & \$225M / \$1.5B & 转型企业AI平台; Socher背景 \\
Kagi        & 全球 & 数十万
  & 未披露 & 纯付费无广告; 极客口碑 \\
Phind       & 全球 & 数百万
  & \$10.9M & 开发者垂直; 可视化; YC \\
\midrule
纳米AI搜索  & 中国 & 9000万+
  & 360集团旗下 & 慢思考模式; 16模型协作 \\
秘塔AI搜索  & 中国 & 6500万
  & 早期融资 & 无广告; 学术; 思维导图 \\
天工AI      & 中国 & 5000万
  & 昆仑万维旗下 & 边搜边创; 自媒体友好 \\
夸克        & 中国 & 3400万DAU
  & 阿里系 & 小而美; 无广告; iOS超百度 \\
开搜        & 中国 & 3000万
  & 早期融资 & 极速响应; 可信度标注 \\
\bottomrule
\end{tabularx}
\end{table}

\subsection{AI搜索的商业模式之争}
\label{subsec:search-business-model}

AI搜索赛道面临一个根本性的商业模式问题：\textbf{如何在"直接给答案"
的模式下赚钱？}

传统搜索引擎的商业模式建立在"用户点击链接"的基础上——Google的收入
来自广告主为搜索结果页上的点击付费（CPC）。但当AI搜索直接给出答案、
用户不再需要点击任何链接时，这一模式就被打破了。

目前AI搜索公司正在探索四种商业模式：

\textbf{模式一：订阅制}。Perplexity Pro（\$20/月）、Kagi Professional
（\$10/月）。优点是收入稳定可预测，缺点是免费替代品（ChatGPT、
Google AI Overviews）的压力使得付费转化率低。Perplexity的数据
表明，订阅仍然是其最大的收入来源，但增长率可能会随着市场饱和
而放缓。

\textbf{模式二：上下文广告}。在AI生成的回答中嵌入赞助内容或
推荐。Perplexity从2025年开始测试这一模式。其优势在于广告可以
与回答内容高度相关（"你问的这个相机，这家店正在打折"），用户
体验损害相对较小。但风险在于——一旦用户感觉到回答的中立性被
广告影响，信任就会崩塌。

\textbf{模式三：企业API和白标}。将AI搜索能力以API形式出售给
企业客户，或为企业构建私有化的AI搜索系统。You.com的转型方向
就是这个。这种B2B模式的好处是客单价高、收入稳定，但需要完全
不同的销售和服务能力。

\textbf{模式四：交易撮合}。当用户搜索商品或服务时，AI直接帮用户
完成购买或预订，搜索引擎从中收取佣金。Perplexity的"in-line
commerce"功能正在探索这一方向。如果成功，这可能是最有想象力的
商业模式——因为它将搜索从"信息发现"推进到了"交易完成"。

对比来看，Google的AI Overviews（在搜索结果页顶部直接展示AI生成的
回答）采用的是"免费+广告"模式——AI回答本身不收费，但Google在AI回答
下方继续展示广告链接。这种模式利用了Google已有的广告基础设施和
客户关系，是AI搜索领域最大规模的商业化实践。

最终，AI搜索的商业模式可能不是四选一，而是四者的混合——就像
YouTube同时拥有广告模式和Premium订阅模式一样。但不同公司的
"混合比例"将决定它们的品牌定位和用户体验——在"用户信任"和
"变现压力"之间找到平衡，是每一家AI搜索公司最大的挑战。

除本节已详细分析的玩家外，AI 搜索赛道还有一批值得关注的独立参与者：

\textbf{Devv.ai}——一个专为开发者设计的 AI 搜索引擎。Devv 的核心
优势是对代码仓库、技术文档和 Stack Overflow 等开发者内容的深度理解。
当开发者搜索一个编程问题时，Devv 不仅给出答案，还能直接展示相关
的代码片段和来自 GitHub 的实现示例。在中国开发者社区中口碑颇高，
被 AIHUB101 等 AI 导航站列为"面向程序员的新一代 AI 搜索引擎"。

\textbf{Globe Explorer}——一款以知识图谱可视化为特色的 AI 搜索工具，
用户输入一个主题后，AI 生成一张交互式的概念地图，展示该主题的
子话题、关联概念和知识层级。这种"视觉化搜索"对于学习新领域和
理解复杂概念特别有用——它不只给你答案，还给你"知识的地图"。

\textbf{Exa.ai}——专门为 AI 模型和开发者设计的搜索 API 平台。
Exa 的核心用户不是终端消费者，而是构建 AI 应用的开发者——他们需要
一个高质量的网络搜索 API 来为自己的 RAG（检索增强生成）管道提供
数据。Exa 在 FRAMES、Tip-of-Tongue 等搜索评测基准中领先，被数千家
公司用于构建 AI 驱动的搜索和检索功能。Exa 代表了 AI 搜索赛道的
"基础设施层"——不面向用户，而是赋能其他 AI 应用。

\textbf{Consensus}——一个专注于学术研究的 AI 搜索引擎，已在
本书第八章"AI 学术工具"一节中详细分析。Consensus 的"共识度量"
功能在 AI 搜索中独树一帜——它不仅告诉你"论文说了什么"，还告诉你
"科学界在这个问题上是否达成了共识"。

\textbf{知乎直答}——知乎推出的 AI 搜索功能，利用知乎平台上数亿条
高质量问答内容作为 AI 的知识基础。知乎直答的独特价值在于其内容
生态——知乎上的回答通常比普通网页更有深度和专业性，这使得知乎直答
在"需要深度解释"的搜索场景中表现出色。这也是"内容平台 + AI 搜索"
模式的典型案例。

\textbf{小红书"点点"}——小红书推出的 AI 搜索应用，深耕生活场景
搜索。当用户搜索"北京周末带娃去哪玩"或"油皮平价防晒推荐"时，
点点从小红书海量笔记中提取答案并以结构化方式呈现。点点代表了
\textbf{垂直内容 + AI 搜索}的独特模式——小红书的内容优势在
"生活经验"领域远超通用搜索引擎。

在 AI 浏览器和搜索的交叉地带，还有一类重要的产品形态——\textbf{AI
浏览器扩展}。这些扩展插件让用户无需切换浏览器就能获得 AI 辅助：

\textbf{Monica AI}——Chrome Web Store 上拥有 300 万+ 用户的 AI
浏览器扩展，由新加坡公司 Butterfly Effect 开发。按下 Cmd/Ctrl+M
即可在任何网页上召唤 AI 助手，支持 GPT-5、Claude、Gemini、DeepSeek
等多个模型。功能涵盖对话、写作、翻译、搜索、图像生成和视频生成，
还内置了 AI Agent 功能。Monica 是目前全球使用量最大的独立 AI 浏览器
扩展之一。

\textbf{Sider AI}——Chrome Web Store 上拥有 500 万+ 用户、11 万+
五星评分的 AI 侧边栏扩展。Sider 同样整合了 GPT-5、Claude、DeepSeek、
Gemini、Grok 等多个模型，提供 AI 搜索、阅读辅助和写作功能。Sider
的评分（4.9 分）和用户量使其成为 AI 浏览器扩展赛道的实际领导者。

\textbf{Merlin AI}——另一款知名的多模型 AI 浏览器扩展，以"一键
总结网页"和"AI 辅助写邮件"功能著称。Merlin 在 Product Hunt 上
多次被推荐，是早期 AI 浏览器扩展赛道的先行者之一。

\textbf{MaxAI}——专注于"阅读增强"的 AI 浏览器扩展，核心功能是
一键总结长文和对网页内容进行 Q\&A。MaxAI 的定位比 Monica 和 Sider
更聚焦——它不试图做"AI 全能助手"，而是专精"阅读效率提升"。

\textbf{豆包浏览器插件}和\textbf{Kimi 浏览器插件}——字节跳动和
月之暗面分别推出的官方浏览器扩展，将其 AI 助手能力带入浏览器
环境。这些官方插件的意义在于打通了"独立 AI App"和"浏览器场景"
——用户无需在两个窗口之间切换，可以在浏览网页的同时调用豆包或
Kimi 的全部能力。

AI 浏览器扩展赛道的竞争格局呈现出明显的"赢家通吃"趋势——Sider
和 Monica 合计占据了大部分市场份额，后来者很难在用户习惯已经形成的
情况下追赶。但这个赛道的最大威胁来自浏览器本身的 AI 化——当 Chrome
内置了 Gemini、Edge 内置了 Copilot 时，独立 AI 扩展的增量价值
正在被压缩。


% =============================================================
\section{AI 浏览器}
\label{sec:social-browsers}
% =============================================================

浏览器是用户接触互联网的第一个界面。当AI被嵌入浏览器时，
它不再是一个需要用户主动打开的工具，而成为了一个\textbf{时刻
在场的助手}——你浏览任何网页时，AI都在旁边，随时准备总结、翻译、
解答和建议。

% -----------------------------------------------------------
\subsection{Arc / Dia $\to$ Atlassian：被收购的AI浏览器梦}
\label{subsec:arc-browser}
% -----------------------------------------------------------

The Browser Company的故事是AI浏览器赛道最戏剧性的篇章。

2022年，由Josh Miller创立的The Browser Company推出了Arc浏览器
——一款从零设计、强调AI功能和工作效率的新型浏览器。Arc以其
激进的界面设计（侧边栏取代传统标签栏、Space分组、自动归档
等功能）迅速在设计师、开发者和科技早期采用者中获得了狂热的
口碑。Arc不只是一个浏览器，它试图\textbf{重新定义人与网页的
交互方式}。公司累计融资约7.47亿美元，员工约93人。

但Arc也面临一个残酷的现实：它吸引了一群热爱折腾的Power User，
却未能突破大众市场。浏览器的网络效应极强——当你的所有书签、密码、
扩展都在Chrome中时，切换成本是巨大的。

2024年底，The Browser Company做出了一个令业界意外的决定：
\textbf{停止Arc的新功能开发}，转而投入开发全新的AI浏览器Dia。
Miller在X上解释了这一决策的逻辑：Arc证明了新型浏览器有市场，
但它的设计对Power User的依赖太重，无法成为"大众浏览器"。Dia的
愿景更加激进——它不仅是一个浏览器，更是一个AI代理，能够理解你
正在做什么，并主动提供帮助：自动帮你填表、总结你正在阅读的文章、
在你购物时比较不同网站的价格。

然而，Dia还未正式发布，命运就已注定。2025年9月4日，Atlassian
宣布以\textbf{6.1亿美元现金}收购The Browser Company。收购逻辑
很清晰：Atlassian的Jira、Confluence、Trello等产品全部运行在浏览器
中，一个AI驱动的"工作浏览器"与Atlassian的生产力套件有着天然的
协同。Atlassian CEO Mike Cannon-Brookes在博客中写道："今天的
浏览器不是为工作设计的。它们是为浏览设计的——看新闻、看视频、
查食谱。但你那些密密麻麻的标签页，每一个都代表一个需要完成的
任务。"

值得一提的是，收购前OpenAI和Perplexity据报道都曾考虑收购
The Browser Company——这说明AI浏览器被各方视为战略性的入口资产。
Miller表示，The Browser Company将在Atlassian旗下独立运营，
继续开发Dia。

6.1亿美元的收购价格相对于7.47亿美元的总融资来说并不算高——这意味着
许多后期投资者可能仅获得了平本或小幅回报。但考虑到Arc从未实现大规模
商业化，这个结果对创始人和早期投资者来说已经算是不错的"软着陆"。

% -----------------------------------------------------------
\subsection{Brave Leo：隐私优先的AI助手}
\label{subsec:brave-leo}
% -----------------------------------------------------------

Brave浏览器以隐私保护著称（默认屏蔽追踪器和广告），其AI助手
\textbf{Leo}于2023年推出，内置于浏览器侧边栏。Leo的设计继承了
Brave的隐私DNA：默认使用本地处理和匿名API调用，不存储用户的
提问数据，不用对话内容训练模型。

用户可以选择不同的AI模型后端，包括Meta的Llama、Anthropic的
Claude和Mistral。免费版提供基础功能，付费版（Leo Premium，
\$6.99/月）提供更高级的模型、更高的使用上限和优先响应。

Leo的功能包括：网页摘要（"总结这个页面"）、内容问答（"这篇文章
的主要论点是什么？"）、翻译、代码解释和创意写作。这些功能与
ChatGPT的能力有重叠，但Leo的独特价值在于它与浏览上下文的
无缝集成——你不需要复制粘贴网页内容到另一个App中。

对于隐私敏感的用户来说——记者、律师、活动人士、政治敏感地区的
用户——Brave Leo提供了一个重要的选择：你可以在浏览敏感内容时
获得AI辅助，而不用担心你的查询被记录、分析或出售。

Brave整体拥有约7000万月活用户，其中Leo的使用率在逐步增长但尚未
成为主流功能。Brave的商业模式围绕其BAT（Basic Attention Token）
加密代币和隐私优先广告展开——Leo的加入为Brave提供了新的付费转化
路径。从战略角度看，Brave Leo代表了一种\textbf{价值观驱动的AI产品}
——隐私不是附加功能，而是产品的核心承诺。

% -----------------------------------------------------------
\subsection{Opera AI 与 Microsoft Edge Copilot}
\label{subsec:opera-edge}
% -----------------------------------------------------------

Opera在2023年推出了"Aria"AI助手，集成在浏览器侧边栏中。Opera的
差异化策略是将AI与其已有的流媒体、游戏和社交功能结合——例如在看
YouTube视频时一键生成摘要，在浏览商品页面时自动比价。Opera还推出了
Opera One浏览器，以"AI原生设计"为卖点。

Opera在AI领域的投入力度和重视程度值得关注。2024年，Opera推出了
Opera One Developer浏览器，将AI作为首要设计原则而非附加功能。
Aria支持联网搜索、图像生成和代码执行等功能，并与Opera的Flow
（跨设备同步工具）和Workspaces（工作空间管理）深度集成。Opera
的活跃用户约3.8亿，虽然大部分集中在非洲和东南亚等新兴市场，但
这些市场对AI工具的需求增速可能比成熟市场更快。

Microsoft Edge Copilot是大厂路线的代表。Edge将Copilot
（基于GPT-4/GPT-4o）直接嵌入浏览器侧边栏，用户可以在浏览任何
网页时召唤AI进行摘要、翻译、问答和创作。Copilot还能直接与Office
365中的文档互动——例如，你可以让Copilot"根据我屏幕上的这篇文章
写一封回复邮件"。Edge Copilot还支持"Compose"功能——用户可以
直接在浏览器中用AI撰写社交媒体帖子、邮件或博客文章，并一键插入
到当前页面。

得益于Microsoft 365的庞大用户基数（超过4亿付费用户）和Windows的
市场份额，Edge Copilot可能是全球使用量最大的浏览器内置AI。
但其体验是否比独立使用ChatGPT更好，仍然见仁见智——许多用户反映
Copilot的侧边栏界面在屏幕空间有限的笔记本电脑上显得拥挤。

AI浏览器赛道的核心问题是：\textbf{浏览器内置AI的价值是否足以让
用户切换浏览器？}对大多数用户来说，Chrome的生态优势（扩展、
同步、密码管理、与Android/Gmail的深度集成）构成了极高的迁移成本。
目前的数据显示，AI功能本身还不足以驱动大规模的浏览器切换——这也是
为什么Atlassian以相对"便宜"的6.1亿美元就收购了The Browser Company，
以及为什么Chrome（72\%市场份额）的地位几乎没有被动摇。

更可能的演进路径是：\textbf{Chrome自己加入AI功能}（Google已经
在Chrome中集成了Gemini Nano的端侧能力和"Help me write"等功能），
从而在不需要用户切换浏览器的情况下完成AI化升级。在这种情况下，
独立AI浏览器的生存空间会被进一步压缩。


% =============================================================
\section{OS 级 AI 集成}
\label{sec:social-os}
% =============================================================

如果说AI浏览器是在"入口"层面嵌入智能，OS级AI则是在更底层的
\textbf{操作系统}层面植入AI。这意味着数十亿设备的用户，无需安装
任何应用，就能直接使用AI功能。这是AI史上最大规模的"默认分发"
——也是对独立AI应用最大的结构性威胁。

% -----------------------------------------------------------
\subsection{Apple Intelligence：苹果的AI觉醒}
\label{subsec:apple-intelligence}
% -----------------------------------------------------------

苹果在AI赛道上一直被外界认为"落后"。当ChatGPT在2022年底引爆
全球时，Siri仍然是一个"只会设闹钟和查天气"的语音助手，在AI圈中
几乎是一个笑话。但2024年6月的WWDC上，苹果发布了
\textbf{Apple Intelligence}，正式加入了这场AI军备竞赛。

Apple Intelligence的核心设计哲学是\textbf{隐私优先}：大部分AI处理
在设备端（on-device）完成，使用苹果自研的芯片（A17 Pro及更新款）
驱动端侧推理；只有复杂任务才通过"Private Cloud Compute"发送到
苹果自己的云端——而这些云端服务器运行的也是苹果自研芯片，且被
设计为不存储用户数据。这一架构意味着用户的数据不会离开苹果的
生态系统，也不会被用于模型训练——这与ChatGPT或Gemini的云端架构
形成了鲜明对比。

Apple Intelligence在iOS 26（2025年秋季发布）中获得了全面部署。
截至2026年2月12日，iOS 26的采用率数据显示：

\begin{itemize}
  \item 66\%的所有iPhone运行iOS 26
  \item 74\%的近四年购买的iPhone运行iOS 26
  \item 57\%的所有iPad运行iPadOS 26
  \item 66\%的近四年iPad运行iPadOS 26
\end{itemize}

这些数字与iOS 18在发布后同期的采用率基本持平（iOS 18在2025年1月：
68\%的所有iPhone），说明Apple Intelligence的存在既没有显著加速
也没有减缓系统升级——用户升级系统的动力仍然是综合性的，而非
单纯为了AI功能。

但66\%的iPhone采用率意味着Apple Intelligence已经覆盖了
\textbf{超过10亿台设备}（iPhone全球活跃安装量约16亿台）。
其核心功能包括：

\begin{itemize}
  \item \textbf{智能写作（Writing Tools）}：系统级的文本重写、
    校对、摘要和语气调整功能，在Mail、Notes、Messages等所有
    文本输入框中可用——不需要安装任何第三方应用
  \item \textbf{Genmoji}：用自然语言描述即可生成自定义Emoji，
    这个功能看似轻巧，但在年轻用户中获得了极高的使用率
  \item \textbf{Image Playground}：在设备端生成个性化图像，
    支持"素描""插画""动画"三种风格
  \item \textbf{增强版Siri}：集成了大语言模型能力，能够理解
    上下文、处理复杂的多步骤指令、跨应用操作。虽然仍不如
    ChatGPT强大，但相比旧版Siri是一个质的飞跃
  \item \textbf{通知摘要}：AI自动对通知进行分类和摘要，让用户
    一眼看到重要信息而不需要逐条阅读
  \item \textbf{Siri + ChatGPT集成}：当Siri无法回答问题时，
    可以选择将请求转发给ChatGPT（用户需主动同意）。这一合作
    对OpenAI来说意义重大——它获得了iPhone用户基础的曝光
  \item \textbf{Clean Up}：照片编辑中的AI橡皮擦功能
  \item \textbf{智能回复}：Mail和Messages中的AI建议回复
\end{itemize}

Apple Intelligence的战略意义不在于某一个功能的惊艳（说实话，
每一个功能单独来看都不如ChatGPT或Claude强大），而在于它将AI
变成了一种\textbf{基础设施}——像Wi-Fi一样，你不需要刻意使用它，
它就在那里，融入了你使用手机的每一个环节。苹果全球约22亿台
活跃设备的安装基数，使得Apple Intelligence从第一天起就拥有了
任何AI创业公司都无法比拟的分发能力。

批评者指出，Apple Intelligence的功能在发布初期相当保守——与
ChatGPT或Claude的能力相比，差距明显。苹果选择了\textbf{安全而非
冒险}：宁可功能少一点、慢一点，也不要出现幻觉和错误引发的PR
危机。2024年12月，Apple Intelligence的通知摘要功能因错误地
"总结"了BBC的新闻标题而引发争议——苹果随后迅速修正了该功能。
这一小事故恰好说明了在数十亿设备上部署AI的风险有多高——任何一个
小bug都会在全球范围内被放大。

% -----------------------------------------------------------
\subsection{Google Gemini Android}
\label{subsec:gemini-android}
% -----------------------------------------------------------

Google在Android端的AI集成走了比苹果更激进的路线。Gemini正在全面
替代Google Assistant，成为Android系统的默认AI助手。

在Pixel手机上，Gemini的集成尤其深入：

\begin{itemize}
  \item \textbf{Gemini Live}：实时语音对话，支持多语言切换，
    可以在锁屏状态下通过"Hey Google"唤醒并使用
  \item \textbf{Gemini Nano}：在设备端运行的轻量级模型（约3B参数），
    用于通话摘要、消息建议、录音机转写等低延迟场景。Gemini Nano
    不需要网络连接即可工作
  \item \textbf{跨应用操作}：Gemini可以理解屏幕内容，在不同应用
    之间传递上下文。例如，你在Google Maps上看一个餐厅页面，可以
    直接对Gemini说"帮我在这家店订今晚的位子"
  \item \textbf{个人智能}：分析Gmail、Drive、Calendar中的内容，
    提供个性化建议——"你下周三有个重要会议，需要我帮你准备汇报
    材料吗？"
  \item \textbf{Circle to Search}：在任何屏幕上画圈即可搜索
    圈中的内容——这是Google与三星合作推出的杀手级功能
\end{itemize}

2026年3月，Google将Gemini的个人智能功能向\textbf{所有美国用户}
开放（此前仅限等候名单），移除了之前的等候限制。这被多家媒体
报道为"有史以来最大规模的消费级AI部署之一"。对Google来说，
Gemini不仅是一个AI助手，更是保护其搜索广告收入的\textbf{防御性
武器}。将Gemini深度嵌入Android和Google Apps，是确保用户留在
Google生态系统内的关键举措。

% -----------------------------------------------------------
\subsection{Samsung Galaxy AI}
\label{subsec:samsung-galaxy-ai}
% -----------------------------------------------------------

三星在2024年1月随Galaxy S24系列推出了\textbf{Galaxy AI}——全球
最早将AI功能作为核心卖点进行大规模营销的智能手机品牌。Galaxy AI
的功能包括：

\begin{itemize}
  \item \textbf{实时通话翻译}：通话中实时翻译对方的语言
  \item \textbf{AI照片编辑}："Generative Edit"可以用AI擦除
    照片中的不需要元素并智能填充
  \item \textbf{Circle to Search}：与Google合作的画圈搜索功能
  \item \textbf{Chat Assist}：在三星消息应用中提供AI建议的
    语气调整和翻译
  \item \textbf{Note Assist}：AI驱动的笔记摘要和格式化
  \item \textbf{Browsing Assist}：浏览网页时的AI辅助摘要
\end{itemize}

三星的策略是与Google深度合作（Galaxy AI底层使用Gemini模型），
同时保持自有品牌。作为全球出货量最大的智能手机厂商之一（2025年
约2.26亿台），三星的AI功能覆盖了庞大的用户基数。Galaxy S25系列
进一步强化了AI功能，将Gemini深度集成为"默认助手"。

值得注意的是，中国手机厂商（小米、OPPO、vivo、荣耀）也在2025年
全面跟进端侧AI功能，但更多依赖自研或国产大模型——例如小米使用了
MiLM和豆包模型，OPPO使用了自研的AndesGPT，vivo使用了BlueLM。
这些厂商的AI功能在中国市场可能比Samsung Galaxy AI更接地气。

% -----------------------------------------------------------
\subsection{Windows Copilot}
\label{subsec:windows-copilot}
% -----------------------------------------------------------

Microsoft将Copilot深度集成在Windows 11中——从任务栏的Copilot按钮
到系统设置中的AI功能，再到与Office 365的无缝联动。Windows Copilot
基于GPT-4/GPT-4o模型，月活约1.05亿（含所有Copilot入口，
FirstPageSage, 2026年3月）。

微软的战略是将Copilot打造为一个\textbf{跨平台的AI"层"}：无论你在
Windows桌面、Edge浏览器、Teams会议、Outlook邮箱还是Word文档中，
Copilot都在那里。2025年推出的\textbf{Copilot+ PC}硬件标准（要求
NPU算力达到45+ TOPS）则从硬件层面为端侧AI奠定了基础。支持Copilot+
PC的笔记本电脑可以运行Recall（AI屏幕记忆）、Live Captions
（实时字幕翻译）等本地AI功能。

OS级AI集成的终极问题是：\textbf{谁拥有AI入口，谁就拥有用户}。
苹果、Google、微软和三星正在将AI变成操作系统的原生能力，这对独立
AI应用（如ChatGPT）构成了长期的结构性威胁。当用户可以直接对着手机
说"帮我总结今天的邮件"而不需要打开任何App时，独立AI助手的使用场景
就被压缩了。

反过来，OpenAI推出ChatGPT桌面客户端和语音模式，Anthropic推出
Claude桌面应用和MCP（Model Context Protocol），都是为了在OS层面
争夺入口——一场静悄悄的平台战争已经开始。这场战争的结果将决定：
在未来的AI时代，\textbf{入口归平台还是归模型公司}？

\subsection{OS级AI的四方对峙格局}
\label{subsec:os-ai-landscape}

四大OS级AI（Apple Intelligence、Google Gemini、Samsung Galaxy AI、
Windows Copilot）之间的竞争可以从四个维度来分析：

\textbf{维度一：端侧能力 vs. 云端能力}。苹果最强调端侧处理
（隐私驱动），大部分AI任务在A17 Pro/M系列芯片上本地运行；
Google走混合路线，Gemini Nano处理简单任务、Gemini Pro/Ultra
处理复杂任务；微软则更依赖Azure云端（Copilot的核心算力在云端），
但通过NPU标准（45+ TOPS）开始建设端侧能力。三星主要依赖Google的
模型能力，自身在端侧AI芯片（Exynos）上的投入相对较弱。

\textbf{维度二：用户覆盖范围}。按设备安装量：Windows约15亿台、
Android约30亿台（其中三星约10亿台）、iOS/macOS约22亿台。从纯
覆盖范围来看，Google/Android拥有最大的潜在触达——但Android的碎片化
意味着许多低端设备无法运行AI功能。苹果的优势在于用户群的\textbf{质量}
——iPhone用户的消费力和应用使用深度平均远高于Android用户。

\textbf{维度三：数据壁垒}。Google拥有最全面的用户数据——搜索历史、
Gmail、YouTube观看记录、Maps定位、Chrome浏览记录。苹果拥有最深度
的设备使用数据——App使用时长、健康数据、支付记录——但受隐私政策
限制，苹果自己也不太"使用"这些数据。微软拥有最丰富的\textbf{工作
数据}——Office文档、Teams会议记录、Outlook邮件。三星的数据资产
相对最弱，主要来自设备传感器和三星健康。

\textbf{维度四：AI模型自主权}。Google完全自主（Gemini是自研模型），
苹果部分自主（自研端侧模型+ChatGPT合作），微软深度绑定OpenAI
（GPT系列），三星完全依赖Google。模型自主权的差异意味着不同的
长期战略灵活性——Google可以根据自己的产品需求自由定制模型行为，
而三星则受制于Google的模型决策。

对于独立AI创业公司来说，OS级AI的崛起既是威胁也是机遇。威胁在于
入口被截流；机遇在于OS级AI的功能往往是"通用而浅层"的——它们无法
像Perplexity那样做深度搜索，无法像Sierra那样理解复杂的业务流程，
无法像Character.ai那样提供深度的角色扮演体验。\textbf{通用入口
会被平台占据，但垂直深度仍然属于创业公司}——这是AI应用层创业
最重要的战略定位原则。

中国市场的OS级AI格局则完全不同。由于Google服务在中国不可用，
Android手机的AI功能由各厂商自行开发：小米使用MiLM和豆包模型，
OPPO使用AndesGPT，vivo使用BlueLM，荣耀与百度合作。这种碎片化
意味着中国市场不存在像Apple Intelligence或Google Gemini那样的
统一OS级AI体验——反而为豆包这样的独立AI应用提供了更大的生存空间。
这也解释了为什么豆包能在中国达到2.27亿月活——在没有统一OS级AI的
市场中，独立应用反而成了"事实上的系统级AI"。

在OS级AI和独立助手之间，还存在一个值得关注的细分赛道——
\textbf{AI硬件助手}。2024年涌现了多款尝试以独立硬件形态承载AI
的设备，但结果大多令人失望：

\textbf{Rabbit R1}——由Jesse Lyu创立的Rabbit公司推出的手掌大小的
AI助手设备，售价\$199，承诺通过"大行动模型"（Large Action Model）
理解用户意图并代理完成任务（如点外卖、打车、搜索信息）。R1在CES
2024上引发轰动，但实际发货后评价极差——核心功能不完善，响应速度慢，
大多数"代理"功能无法可靠工作。据报道仅售出约15,000台后即停产。
Rabbit R1的失败是"AI硬件泡沫"的标志性案例——在AI软件仍在快速
进化的阶段，将AI锁定在一个固定硬件形态中可能为时过早。

\textbf{Humane AI Pin}——由前苹果高管Imran Chaudhri和Bethany
Bongiorno创立，这款激光投影式的AI别针设备承诺"后智能手机时代"
的交互方式。AI Pin售价\$699+月费\$24，可以将信息投射到手掌上、
通过语音进行AI交互。但产品发布后遭遇了几乎一致的负面评价——
过热、电池续航差、投影效果差、AI响应不准确。Humane在2024年10月
暂停销售并裁员，是2024年最惨烈的AI硬件失败案例之一。

\textbf{Tab}（Avi Schiffmann）——18岁的天才少年Avi Schiffmann
开发的AI可穿戴项链，试图通过全天候被动聆听和记录来构建"AI记忆"。
Tab的理念是不需要用户主动提问——AI始终在听，在合适的时机提供
相关信息和建议。虽然引发了极大的隐私争议，但Tab代表了AI硬件
最前沿的思考——从"用户发起交互"到"AI主动感知环境"。

与这些"后智能手机"实验相比，中国市场的AI硬件走了一条更务实的
路线——将AI嵌入已有的成熟硬件品类中：

\textbf{小度}（百度旗下）——中国最大的智能音箱品牌之一，已将
文心一言大模型整合进小度系列产品。用户可以与小度进行多轮对话、
让AI讲故事、辅导作业，甚至进行简单的情感陪伴。小度的策略是
"AI让现有硬件更智能"而非"创造全新的AI硬件品类"。

\textbf{天猫精灵}（阿里巴巴旗下）——同样将通义千问大模型接入了
天猫精灵智能音箱，升级了对话能力和任务执行能力。天猫精灵还深度
整合了阿里电商生态，用户可以通过语音AI直接在淘宝/天猫下单。

中国AI硬件赛道与美国市场的最大差异是：中国厂商更倾向于在\textbf{
已有硬件形态}（音箱、学习机、手机）上叠加AI能力，而非创造全新的
硬件品类。Rabbit R1和Humane AI Pin的失败在某种程度上验证了中国
厂商的务实策略——至少在当前阶段，"AI增强已有设备"比"AI驱动
全新设备"更具商业可行性。


% =============================================================
\section{AI 销售与营销工具}
\label{sec:social-sales}
% =============================================================

从本节开始，我们从纯C端场景过渡到\textbf{前台业务场景}——销售、
客服和HR。这三个领域的共同特征是：大量重复性的人际沟通、丰富的
交互数据、以及对效率提升的强烈需求。AI在这些场景中正在从"分析工具"
进化为"行动代理"。

从市场规模来看，仅AI销售工具赛道就已经达到数百亿美元的TAM
（Total Addressable Market）。全球有超过3亿名从事销售和营销
相关工作的专业人员，如果每人每年的AI工具支出达到1000美元，
就是一个3000亿美元的市场。当然，这个理论上限还很遥远——但Gong
以超过3亿美元ARR证明了企业愿意为AI销售工具付出真金白银，Clay
的6倍估值跳升证明了资本市场对这一赛道的热情。

% -----------------------------------------------------------
\subsection{Gong：对话智能的先驱}
\label{subsec:gong}
% -----------------------------------------------------------

Gong是AI销售工具赛道的标杆公司。成立于2016年，总部位于旧金山，
Gong的核心产品是一个\textbf{收入智能平台}（Revenue Intelligence
Platform）——它自动录制和分析销售团队与客户的通话、邮件和会议，
从中提取关于交易健康度、客户情绪、竞品提及、关键决策者参与度等
多维度的洞察。

Gong在2021年完成了由Franklin Templeton领投的2.5亿美元E轮融资，
估值72.5亿美元，总融资5.84亿美元。Coatue、Salesforce Ventures、
Sequoia、Thrive Capital和Tiger Global等顶级投资者均参与了投资。

2025年3月，Gong宣布ARR突破\textbf{3亿美元}——对于一家自2021年以来
没有进行后续融资的公司来说，这一增速表明其商业模式已经充分验证。
公司拥有超过1000名员工，客户覆盖已从纯科技行业扩展到金融、医疗、
制造等更广泛的垂直领域。

生成式AI的加入让Gong的产品能力实现了质的飞跃。公司的"Ask Anything"
功能让销售经理可以用自然语言查询所有通话和交易数据——"哪些交易
在上个月提到了竞品X？""这个季度我们在医疗行业的成单率比上季度
高还是低？""帮我总结这个客户过去三次通话的核心关切。"这种
自然语言查询能力在ChatGPT之前是不可想象的——之前的做法是通过
预设的Dashboard和报表来看数据，维度固定且不灵活。

Gong的AI能力使用量在过去12个月内增长了近50\%。更重要的是，Gong
正在从分析工具向\textbf{行动工具}转型——不仅告诉你发生了什么，
还主动推荐下一步行动（"这个交易的风险信号增加了，建议在本周安排
一次与VP的跟进电话"）。这一趋势与整个AI赛道从"辅助"到"代理"
（agentic）的演进方向一致。

从IPO角度看，超过3亿美元的ARR加上持续增长，使得Gong已经具备了
上市条件。多家投行已在讨论Gong的IPO时间线，如果成功，它将成为
AI销售工具赛道的旗舰上市公司。

% -----------------------------------------------------------
\subsection{Clay：增长最快的AI销售自动化}
\label{subsec:clay}
% -----------------------------------------------------------

Clay是AI销售工具赛道中增长最快、最令投资者兴奋的黑马。

成立于2017年，Clay的核心产品是一个\textbf{AI驱动的GTM
（Go-To-Market）平台}，帮助销售和营销团队进行数据富集
（data enrichment）、潜客发掘（lead discovery）和个性化
外联（personalized outreach）。

Clay的独特之处在于其\textbf{数据聚合能力}：它可以从\textbf{150多个
数据源}中自动拉取和交叉验证潜在客户的信息——包括LinkedIn、
公司官网、新闻报道、职位变动、融资事件、技术栈变化等。然后使用
AI对每个潜客生成个性化的外联文案——不是千篇一律的"Dear Sir/Madam"，
而是"我注意到你们上个月刚从AWS迁移到GCP，这个过程中数据安全通常是
最大的挑战……"

客户名单包括OpenAI、Anthropic、Canva、Intercom和Rippling等
知名科技公司——这些公司自己就是AI的深度使用者，他们选择Clay说明
了产品的实力。

2025年是Clay的融资爆发之年：

\begin{itemize}
  \item 2025年2月：12.5亿美元B轮融资
  \item 2025年6月：15亿美元Sequoia领投的老股交易（tender offer），
    允许大部分员工出售部分股份
  \item 2025年8月：1亿美元C轮，估值\textbf{31亿美元}，CapitalG
    （Google的投资基金）领投，Meritech、Sequoia、First Round、
    BoxGroup、Boldstart和Sapphire Ventures跟投
\end{itemize}

总融资2.04亿美元，估值31亿美元——不到一年内估值从5亿美元飙升至
31亿美元，增长超过6倍。更令人印象深刻的是，这家公司已经8岁了
（2017年成立），在AI浪潮到来之前默默耕耘了多年，直到生成式AI的
爆发让它的产品价值被充分释放。

Clay的成功反映了一个重要趋势：\textbf{AI正在重塑整个销售工作流}。
传统的CRM（如Salesforce）是"记录系统"（system of record）——
它记录你做了什么。Clay则试图成为"行动系统"（system of action）
——它告诉你应该做什么，甚至帮你直接做。这是从被动记录到主动行动的
范式转变。

% -----------------------------------------------------------
\subsection{其他重要玩家}
\label{subsec:sales-others}
% -----------------------------------------------------------

\textbf{Apollo.io}是Clay最直接的竞争对手，提供B2B数据库和销售
自动化功能。Apollo的优势在于其\textbf{2.75亿+联系人}的数据库规模
和更低的入门价格（有免费方案），适合中小型销售团队。但在AI能力的
深度和数据富集的广度上弱于Clay。Apollo与Clay的竞争本质上是
"广度vs.深度"的经典商业战略之争。

\textbf{Outreach}是销售参与（Sales Engagement）平台的领导者，
帮助销售团队管理多渠道的客户沟通（邮件、电话、LinkedIn、短信）。
Outreach在2021年估值达到44亿美元，并在2024--2025年大幅加强了
AI功能，包括AI驱动的销售预测（pipeline forecasting）和自动化
销售工作流。Outreach与Gong的差异在于：Gong侧重于"对话分析"
（通话后洞察），Outreach侧重于"销售执行"（通话前和通话中的工作流）。

\textbf{Chorus}（ZoomInfo旗下）是Gong在对话智能领域的主要竞品。
2021年被ZoomInfo以5.75亿美元收购后，Chorus的独立品牌逐渐被整合
到ZoomInfo的更广泛平台中。这一整合趋势在AI销售工具赛道中非常
典型——独立的"点状工具"（point solution）越来越难独立生存，
要么被更大的平台收购，要么自己演化为平台。

\begin{keybox}[AI销售与营销工具赛道全景]
\small
\begin{tabularx}{\textwidth}{l c r r X}
\toprule
\rowcolor{figLightGray}
\textbf{公司} & \textbf{成立} & \textbf{融资}
  & \textbf{估值/收入} & \textbf{核心定位} \\
\midrule
Gong        & 2016 & \$584M
  & \$7.25B / \$300M ARR & 对话智能; Ask Anything; 准IPO \\
Clay        & 2017 & \$204M
  & \$3.1B & 数据富集; 150+源; AI外联 \\
Apollo.io   & 2015 & \$244M
  & $\sim$\$1.6B & 2.75亿联系人库; 中小团队友好 \\
Outreach    & 2014 & \$489M
  & \$4.4B & 多渠道销售参与; AI预测 \\
Chorus      & 2015 & 被ZoomInfo收购
  & \$575M收购价 & 对话分析; 整合入ZoomInfo \\
\bottomrule
\end{tabularx}

\medskip
\noindent AI销售工具赛道的下一个大趋势是\textbf{从"辅助"到"代理"}：
AI不再只是帮销售人员分析数据和生成文案，而是开始自主执行完整的
销售动作——自动发送个性化邮件、自动安排会议、自动更新CRM记录。
这一趋势对Salesforce等传统CRM巨头构成了直接威胁，也为Clay等
AI原生公司提供了重新定义"销售软件"品类的窗口。
\end{keybox}

AI销售工具赛道还有一个值得关注的结构性趋势：\textbf{工具整合}。
销售团队目前需要同时使用CRM（Salesforce/HubSpot）+ 数据工具
（Clay/Apollo）+ 参与工具（Outreach/Salesloft）+ 对话分析工具
（Gong/Chorus）+ 教练工具（Second Nature/Rehearsal）。平均一个
销售团队使用超过10个软件工具。这种"工具爆炸"不仅增加了成本，
也降低了效率——销售人员花在管理工具上的时间几乎和花在销售上的
时间一样多。

AI的到来正在加速这一整合进程。Gong正在向完整的"收入平台"演进，
Clay正在成为"GTM的操作系统"，Salesforce则通过Einstein GPT和
Agentforce试图在CRM中集成所有AI能力。未来2--3年，AI销售工具赛道
很可能会经历一轮整合洗牌——独立的点状工具要么被收购整合，要么
自我进化为平台。

在B2B销售工具之外，面向\textbf{营销创意和广告生成}的AI工具也是
一个快速增长的细分赛道：

\textbf{Anyword}——一个专为市场营销文案打造的AI平台，核心差异化
在于其\textbf{预测性能评分}——AI不仅生成广告文案，还能预测每条
文案的点击率和转化率（基于对真实营销数据的A/B测试训练）。Anyword
声称能实现30\%的文案效果提升。平台支持Google Ads、Facebook Ads、
邮件营销等多个渠道的文案生成。对于每天需要批量生产广告变体的
营销团队来说，Anyword解决的核心痛点是"哪条文案会赢"——这是
传统AI写作工具（如Jasper）不具备的能力。

\textbf{Predis.ai}——一个AI广告和社交媒体内容生成器，30秒内即可
创建广告图片、UGC视频和社交媒体帖子。Predis 被全球30万+用户使用，
在多个评测中获得4.8--4.9的高评分。其核心卖点是"从创意到发布的
一站式"——不仅生成内容，还内置排程功能，可以直接发布到各社交
平台。

\textbf{AdCreative.ai}——专注于广告创意生成的AI工具，能在数秒内
生成大量广告图片和文案变体，并按预测的转化率排序。AdCreative 的
独特之处在于其庞大的广告效果数据库——平台根据历史上哪些视觉风格
和文案策略在特定行业中表现最好来指导AI的创意生成。

\textbf{Lately AI}——一个AI内容营销平台，专注于将长内容（博客、
播客、视频）自动拆分为数十条社交媒体帖子。Lately 的AI分析品牌
的历史社交数据，学习哪些风格的帖子在该品牌的受众中表现最好，
然后据此生成新内容。这种"学习品牌语调"的能力使Lately在企业
品牌营销团队中获得了良好口碑。

在中国市场，\textbf{巨量引擎}（字节跳动旗下）是最大的AI广告
平台——其"巨量千川"和"巨量创意"工具利用AI自动生成广告素材、
优化投放策略和预测转化效果。巨量引擎服务于数百万在抖音上投放
广告的商家，是中国短视频广告AI化的最前沿。\textbf{百度营销AI}
则在搜索广告领域提供AI创意生成和智能出价工具，利用文心大模型
为广告主自动撰写搜索广告文案并优化关键词策略。


% =============================================================
\section{AI 客服}
\label{sec:social-customer-service}
% =============================================================

客户服务是AI最早也是最成功的商业应用场景之一。原因很简单：客服场景
有大量重复性问题（40--70\%的客服咨询是重复的）、丰富的历史数据
（每次客户交互都有记录）、明确的成功指标（首次解决率、平均处理时间、
客户满意度CSAT），以及巨大的成本优化空间（一个人类客服的年均成本
在美国约5--8万美元）。

AI客服自动化市场在2026年的规模约为100--110亿美元，增长率23--26\%。
多家行业领导者预测，\textbf{2026年将是AI代理处理的客服交互量首次
超过人类}的元年——这是一个具有标志性意义的转折点。

这一赛道的技术演进可以分为三代：

\textbf{第一代（2015--2020）：规则引擎和决策树}。传统聊天机器人
按照预设的IF-THEN规则运行，只能处理最简单的问题（"营业时间是什么？"）。
用户体验差，解决率低（通常低于20\%），更多是"假装有AI"。

\textbf{第二代（2020--2023）：NLU增强的对话系统}。使用自然语言
理解（NLU）技术提升了意图识别的准确性，但仍然依赖预定义的意图分类
和对话流程。解决率提升到30--50\%，但面对复杂或未预见的问题仍然力不从心。

\textbf{第三代（2023至今）：LLM驱动的自主AI代理}。以Sierra、
Intercom Fin、Zendesk+Forethought为代表。基于大语言模型，能够理解
自由文本、跨系统查询信息、自主决策和执行操作。声称解决率达到
70\%以上。这一代产品才是真正意义上的"AI客服"。

% -----------------------------------------------------------
\subsection{Sierra AI：Bret Taylor 的百亿美金赌注}
\label{subsec:sierra}
% -----------------------------------------------------------

Sierra AI是AI客服赛道中估值最高、最受关注的创业公司——也是整个
AI应用层中为数不多的"十亿美金俱乐部"成员。

Sierra由\textbf{Bret Taylor}和Clay Bavor于2024年初在旧金山联合
创立。Taylor的履历几乎无人能及——他不是一个普通的连续创业者，
而是过去二十年硅谷最重要的技术领袖之一：

\begin{itemize}
  \item Google Maps联合创始人（2004年）
  \item FriendFeed创始人，后被Facebook收购（2009年）
  \item Salesforce CTO，后成为联合CEO（2017--2022年）
  \item Twitter董事会主席（2022年）
  \item \textbf{OpenAI董事会主席}（2024年至今）
\end{itemize}

Bavor则是Google的长期高管（超过18年），曾负责VR/AR部门和多个
实验性项目（包括Google Lens的早期版本）。

Sierra的核心产品是\textbf{自主AI客服代理}（autonomous AI agent）
——这不是传统的聊天机器人（chatbot），后者按照预设的决策树和脚本
运行，遇到超出脚本的问题就卡住。Sierra的AI代理能够理解复杂的
客户问题、跨系统查询信息（CRM、订单管理、支付系统）、自主做出
决策（是否退款、如何升级服务）并执行操作。

客户包括SoFi、Ramp、Brex、WeightWatchers等知名企业。在这些客户
的实际部署中，Sierra的AI代理声称能够自主解决70\%以上的客户咨询，
剩余的30\%自动升级给人类客服并附带完整的上下文。

Sierra在短短18个月内就达到了惊人的规模：

\begin{itemize}
  \item 2024年2月：A轮，Sequoia领投
  \item 2025年1月：追加融资，估值约40亿美元
  \item 2025年9月：3.5亿美元融资，估值\textbf{100亿美元}，
    Greenoaks Capital领投
  \item 总融资：6.35亿美元
  \item 2025年收入：约1.5亿美元
  \item 客户数量：数百家
\end{itemize}

从创立到100亿美元估值仅用了约18个月——Sierra成为了AI赛道中少数
几家与OpenAI、Anthropic、xAI、Safe Superintelligence等基础模型
公司并肩的应用层"超级估值"公司。

Sierra的成功印证了一个重要命题：\textbf{在AI应用层，创始人的行业
信誉、客户网络和信任度是最重要的护城河}。Taylor在Salesforce积累的
数千个企业客户关系，使得Sierra在推出产品后几乎立即获得了顶级客户
的信任。对于企业来说，将客服系统——一个直接面向客户、影响品牌
声誉的关键环节——交给一家初创公司，需要极高的信任度。而Taylor的
名字和履历就是这份信任的最好背书。

% -----------------------------------------------------------
\subsection{Intercom Fin：从客服工具到AI代理}
\label{subsec:intercom-fin}
% -----------------------------------------------------------

Intercom是一家成立于2011年的老牌客户沟通平台，总部位于旧金山和
都柏林。2023年，Intercom推出了AI客服代理\textbf{Fin}，标志着公司
从传统客服工具向AI原生平台的全面转型。

Fin经历了快速迭代。2025年10月，Intercom在其年度Pioneer大会上发布了
\textbf{Fin 3}——宣称是"最强大的AI客服代理"，能够处理跨渠道
（聊天、邮件、电话、社交媒体）的复杂查询。Intercom的CEO Eoghan
McCabe分享了公司的愿景：将Fin从一个"AI聊天机器人"进化为一个
\textbf{统一客户代理}（unified Customer Agent）——不仅回答问题，
还能执行退款、修改订单、查询账户状态、升级服务计划等操作。

Fin的定价模式独特且大胆：按\textbf{解决量}收费（\$0.99/每次成功
解决），而非传统SaaS的按座位数收费。这种模式意味着客户只为实际
价值付费——如果Fin没有解决客户的问题，Intercom不收钱。但在高流量
场景下，\$0.99/次的费用会快速累积，总成本可能显著高于传统方案。

McCabe在Pioneer 2025上做出了一个大胆预测：\textbf{2026年将是
AI代理处理的服务交互首次超过人类客服的一年}。这一预测与Zendesk
的判断不谋而合——两家行业领导者的共识，使得这一预测具有了更高的
可信度。

% -----------------------------------------------------------
\subsection{Zendesk AI 与 Forethought 收购}
\label{subsec:zendesk-forethought}
% -----------------------------------------------------------

Zendesk作为客服SaaS的老牌巨头（2007年成立，曾在纽交所上市，
2022年以105亿美元被私有化），一直在积极拥抱AI。

2026年3月11日，Zendesk宣布进入收购AI客服初创公司
\textbf{Forethought}的最终协议——后者是2018年TechCrunch Battlefield
的冠军，比ChatGPT早了四年就开始做AI客服自动化。Forethought总融资
1.15亿美元，投资者包括Blue Cloud Ventures等。到2025年，Forethought
已支持每月超过\textbf{10亿次}客户交互，客户包括Upwork、Grammarly、
Airtable和Datadog。

Zendesk的收购逻辑很明确：通过Forethought的\textbf{自学习AI代理}
（self-improving AI agents）技术，增强其Resolution Platform的
自主服务能力。交易预计在2026年3月底前完成，收购金额未披露。

Zendesk将这次收购放在一个宏大的叙事中："2026年将是AI代理处理的
服务交互首次超过人类的一年"——这与Intercom的判断完全一致。两家
巨头不约而同地选择在这个时间点加速AI投入，说明AI客服的拐点
确实已经来临。

值得一提的是Forethought的历程本身就是一个引人入胜的创业故事。
公司由Deon Nicholas于2017年创立，比ChatGPT早了五年就开始用AI
解决客服问题。在2018年的TechCrunch Battlefield上，Forethought
在众多参赛创业公司中脱颖而出获得冠军——彼时AI还远未成为主流话题。
此后的七年中，Forethought默默耕耘，建立了扎实的技术栈和客户基础。
当2023年生成式AI浪潮到来时，Forethought已经拥有了数年的客服数据
积累和领域知识——这使得它在引入LLM能力时比从零开始的竞争对手更
具优势。

Forethought的"自学习"特性是其核心差异化：AI代理会从每一次客户
交互中学习，逐步提升解决率和回答质量，而不需要人工干预或频繁的
模型微调。这种"越用越好"的飞轮效应使得客户一旦部署就很难切换
——切换成本不仅包括技术迁移，还包括损失AI积累的领域知识。

Zendesk以非公开金额收购Forethought，结束了后者的独立运营。对于
AI客服创业公司来说，Forethought的结局既是一个成功案例（成功退出），
也是一个警示——在大厂全面拥抱AI的背景下，独立AI客服公司的独立
生存空间越来越小。

% -----------------------------------------------------------
\subsection{Ada：加拿大的AI客服独角兽}
\label{subsec:ada-cx}
% -----------------------------------------------------------

Ada成立于2016年，总部位于多伦多，是最早的AI客服自动化平台之一。
公司将自己定位为"面向企业的AI客服转型伙伴"——不仅提供技术产品，
还提供从规划到上线的全流程咨询服务。

Ada声称其AI代理能够\textbf{解决超过70\%}的客户服务咨询。总融资约
1.91亿美元，最新一轮为2021年的C轮。2026年的数据显示，Ada拥有
约260名员工（同比增长63.5\%），年收入约3100万美元，服务覆盖
19个国家。客户包括Square、YETI和Monday.com。

与Sierra相比，Ada更侧重于\textbf{高流量、标准化}的客服场景
（如电商的"我的订单在哪里？""如何退货？"），而Sierra则瞄准
更复杂的企业级需求（如金融科技的账户问题、保险的理赔流程）。
两者在市场定位上有明确的差异化。

\begin{table}[htbp]
\centering
\caption{AI客服赛道核心公司对比（截至2026年3月）}
\label{tab:ai-cs-comparison}
\small
\begin{tabularx}{\textwidth}{l c r X}
\toprule
\rowcolor{figLightGray}
\textbf{公司} & \textbf{成立} & \textbf{融资/估值}
  & \textbf{核心特色} \\
\midrule
Sierra AI   & 2024 & \$635M / \$10B
  & 自主AI代理; Bret Taylor; 企业级 \\
Intercom Fin & 2011$^*$ & Intercom整体
  & 统一客户代理; \$0.99/解决; 跨渠道 \\
Zendesk AI  & 2007$^*$ & 私有（\$10.2B收购）
  & Resolution平台; 收购Forethought \\
Ada         & 2016 & \$191M
  & 70\%+解决率; 多伦多; 19国覆盖 \\
Forethought & 2018 & \$115M $\to$ 被Zendesk收购
  & 自学习AI; 10亿+月交互; TC冠军 \\
\bottomrule
\end{tabularx}

\medskip
\noindent {\scriptsize $^*$ Intercom和Zendesk的成立年份为公司
整体，非AI产品发布时间。Sierra是唯一的AI原生创业公司。}
\end{table}

AI客服赛道正在经历一场深刻的商业模式变革。传统客服SaaS的商业
模式是按"座位"（agent seat）收费——公司有多少客服人员，就付多少
钱。当AI代理开始替代人类客服时，这一模式面临根本性挑战——如果
你的产品目标是\textbf{减少客服人数}，你怎么能按客服人数收费？

Sierra和Intercom都在探索新的定价模式：Sierra按"AI代理解决的交互量"
收费，Intercom按"每次解决\$0.99"收费。这代表了\textbf{从"为人力
付费"到"为结果付费"}的转变，可能是SaaS行业自云化以来最大的
商业模式创新。

这一转变对整个客服行业的影响是深远的：

\begin{itemize}
  \item 客服外包公司（如Teleperformance、Concentrix）面临
    结构性威胁——如果AI能处理70\%+的咨询，外包座席的需求
    将大幅下降
  \item 客服软件公司（如Zendesk、Freshdesk）需要从"帮助人类
    更高效地工作"转型为"直接替代人类工作"——这要求完全不同的
    产品架构
  \item 企业客户的采购逻辑从"我需要多少个客服"变成"我需要
    多少次高质量的客户交互"
\end{itemize}

在上述主要玩家之外，AI客服赛道还有多个值得关注的独立参与者：

\textbf{Tidio}——一款面向中小企业的AI客服聊天机器人，其Lyro AI
Agent声称拥有业界最高的67\%自动解决率。Tidio的核心优势是"5分钟
部署"——小型电商和SaaS公司无需技术团队即可快速上线AI客服。支持
13种语言，每次AI对话成本仅\$0.50。Tidio的定位精准瞄准了Sierra和
Intercom不太覆盖的中小企业市场——这些企业没有大型客服团队，但
同样需要7×24小时的客户响应能力。

\textbf{Freshdesk Freddy}——Freshworks旗下的AI客服产品线。Freddy
AI集成在Freshdesk客服平台中，提供预构建的代理工作流、AI辅助工单
处理和性能洞察。Freddy的优势在于Freshworks已有的SMB客户基础——
对于已经使用Freshdesk的企业来说，开启AI功能几乎是零迁移成本。

\textbf{Drift}（现为Salesloft旗下）——一款对话式营销和销售AI平台，
2024年被Salesloft收购并整合。Drift的AI聊天机器人专注于B2B网站的
"访客转化"场景——当潜在客户访问你的网站时，AI主动发起对话、
资质审核、安排与销售代表的会议。Drift的收购是AI客服/销售赛道
工具整合趋势的又一案例。

在中国市场，AI客服赛道同样竞争激烈：

\textbf{智齿科技}——中国领先的AI客服SaaS提供商，产品覆盖在线客服、
呼叫中心、工单系统和AI机器人。智齿在2024--2025年间深度整合了大模型
能力，支持接入DeepSeek、GPT等多个模型。核心客户包括多家银行、保险
公司和大型电商平台。

\textbf{网易七鱼}——网易旗下的智能客服平台，以全渠道接入和大模型
驱动的Agent协同为卖点。七鱼的差异化在于其与网易生态（网易邮箱、
网易严选等）的深度整合，以及在AIGC内容审核和情感分析方面的
技术积累。在2025年阿里云开发者社区的深度评测中，七鱼被评为
"AI能力表现突出"的智能客服产品之一。

\textbf{腾讯企点}——腾讯旗下的企业级智能客服平台，深度整合微信
生态。企点AI的核心优势是能够在微信公众号、小程序、企业微信等
腾讯系触点上提供无缝的AI客服体验。对于以微信为主要客户触点的
中国企业来说，企点的生态优势几乎不可替代。


% =============================================================
\section{AI HR 与招聘}
\label{sec:social-hr}
% =============================================================

人力资源和招聘是另一个被AI深度渗透的领域。从简历筛选到面试评估，
从技能匹配到员工留存预测，AI正在重塑HR的每一个环节。Gartner报告
指出，2026年有\textbf{65\%的企业}优先采用"技能导向型招聘"
（skills-based hiring）——这一趋势的背后，是AI驱动的技能匹配
和评估技术的成熟。

% -----------------------------------------------------------
\subsection{HireVue：AI面试的先驱（与争议）}
\label{subsec:hirevue}
% -----------------------------------------------------------

HireVue是最早将AI引入招聘面试的公司之一，成立于2004年，总部位于
犹他州南乔丹。其核心产品是\textbf{AI视频面试平台}——候选人录制
视频回答预设问题，AI分析其语言内容、表达方式和行为信号，给出
结构化的评估分数。

HireVue在2019年的估值一度超过10亿美元，被联合利华、高盛、希尔顿
等大型企业广泛采用。据称已完成超过7000万次AI面试。但其早期版本中
的"AI面部分析"功能引发了巨大争议——批评者（包括电子前哨基金会EFF
和多个学术研究团队）认为通过面部表情和微表情判断候选人的能力和
品格存在\textbf{系统性偏见}，尤其可能对不同种族、文化背景和
神经多样性（neurodivergent）的候选人造成不公平。

2021年，HireVue在持续的外部压力下移除了面部分析功能，转而聚焦于
\textbf{语言内容和结构化评估}。到2025--2026年，HireVue已经完成了
品牌重塑，将自己定位为"AI驱动的结构化面试平台"，强调其评估方法
基于经过验证的工业心理学（Industrial-Organizational Psychology）
框架。

新一代HireVue使用生成式AI进行更细粒度的能力评估，包括：AI辅助的
面试问题生成（根据职位需求自动定制面试题）、实时的候选人表现
评分、以及面试后的AI反馈报告。

HireVue的故事是AI HR赛道的一面镜子：它展示了AI在HR中应用的
\textbf{巨大潜力}（每年处理数千万次面试、大幅提升招聘效率），
也暴露了\textbf{公平性和透明度}方面的严峻挑战。

更广泛地看，HireVue的经历预示了AI在所有"高利害关系"场景中面临的
共同困境：当AI的决策直接影响人的生活（就业、贷款、保释、医疗诊断）
时，社会对公平性和透明度的要求会比对普通工具高出几个量级。AI招聘
工具的开发者需要投入大量资源进行偏见测试、可解释性设计和监管合规
——这些"非功能性"投入不直接创造收入，但却是产品能否被企业采纳的
前提条件。对于创业公司来说，这意味着AI HR赛道的准入门槛不仅是
技术门槛，更是\textbf{信任和合规门槛}。

% -----------------------------------------------------------
\subsection{Paradox $\to$ Workday：AI招聘的最大并购}
\label{subsec:paradox}
% -----------------------------------------------------------

Paradox是AI招聘赛道中最成功的创业公司之一。其核心产品是一个名为
\textbf{Olivia}的对话式AI招聘代理——它可以自动化整个招聘前端
流程：通过短信和聊天接收申请、筛选简历、安排面试时间、回答
候选人关于职位和公司的问题。

Paradox的强项在于\textbf{高流量前线招聘}（frontline hiring）
——零售、餐饮、酒店、物流等行业每年需要招聘数十万甚至数百万名
小时工。传统的人工筛选和安排流程效率极低（从申请到面试通常需要
3--5天），而Olivia将这一过程缩短到\textbf{数分钟}——候选人提交
申请后，AI立即开始对话，确认基本资格，如果匹配就当场安排面试。

2025年8月21日，Workday宣布签署最终收购协议，收购Paradox。
收购金额未披露，但考虑到Paradox是该赛道最大的独立玩家，估值
可能在数亿到十亿美元级别。

收购逻辑非常清晰：Workday是全球最大的HR云平台之一（市值约
750亿美元），但其在AI驱动的人才获取（talent acquisition）
方面存在能力空白。Paradox的AI招聘代理能够填补这一空白，
特别是在Workday相对薄弱的高流量前线招聘场景中。

这笔收购标志着\textbf{AI HR创业公司被大厂整合}的趋势加速。
在Workday、SAP SuccessFactors、Oracle HCM和ADP等HR平台巨头
纷纷加强AI能力的背景下，独立AI HR创业公司的生存窗口正在快速
关闭——如果不能在2--3年内建立足够的规模和客户粘性，就会面临
被收购或边缘化的命运。

% -----------------------------------------------------------
\subsection{Eightfold AI：技能智能的领导者}
\label{subsec:eightfold}
% -----------------------------------------------------------

Eightfold AI成立于2016年，总部位于加州圣克拉拉，总融资3.97亿美元，
估值约21亿美元（2021年E轮，2.2亿美元）。

Eightfold的核心差异化在于其\textbf{人才智能平台}（Talent
Intelligence Platform）——不是简单地用关键词匹配简历和职位描述，
而是用深度学习理解候选人和职位之间的\textbf{深层技能关联}。
例如，一个有5年Python经验的数据分析师可能非常适合一个要求
"机器学习工程师"的职位，即使简历上没有出现"机器学习"这个
关键词——因为AI理解Python和数据分析是机器学习的基础技能。

Eightfold的平台覆盖三个维度：

\begin{itemize}
  \item \textbf{人才获取}（Talent Acquisition）：从外部候选人池中
    识别最佳匹配
  \item \textbf{人才管理}（Talent Management）：帮助现有员工找到
    内部晋升和横向发展机会
  \item \textbf{人才洞察}（Talent Insights）：劳动力市场分析、
    技能缺口识别、多元化与包容性（D\&I）指标追踪
\end{itemize}

但Eightfold面临的挑战是其\textbf{最新融资已过去五年}。在AI赛道
快速迭代的背景下，2021年的估值和技术栈是否仍然具有竞争力，是一个
需要验证的问题。生成式AI的爆发使得HR AI赛道的技术门槛发生了根本
性变化——2021年的"AI"和2025年的"AI"已经不是一个东西了。

Eightfold的创始人Ashutosh Garg（前Google搜索工程师）和Varun
Kacholia（前Facebook新闻推送负责人）拥有深厚的AI和搜索引擎背景，
这使得Eightfold在技能图谱的构建方面具有独特优势。公司声称已经
构建了全球最大的人才技能图谱，覆盖超过10亿个人才档案。这个图谱
不断通过新数据自我更新——每当一个用户在LinkedIn上更新职位、一家
公司发布新职位描述、一篇学术论文提出新技能分类时，图谱都会相应
调整。

对于大型企业（如联合利华、德勤、美国银行等Eightfold的客户）来说，
Eightfold解决的核心问题是\textbf{内部人才的可见性}。在一家拥有
10万+员工的跨国公司中，HR部门往往不知道自己的员工拥有哪些技能、
可以胜任哪些内部岗位。Eightfold的平台让"内部招聘"变得和"外部招聘"
一样高效——这不仅降低了招聘成本，也提高了员工留存率（因为员工
可以在公司内部找到职业发展路径，而不需要离职跳槽）。

Eightfold在多元化和包容性（D\&I）方面的功能也是其卖点之一。平台
可以在筛选过程中屏蔽候选人的姓名、照片和教育背景，只展示技能
匹配度——这种"盲审"机制旨在减少无意识偏见。虽然效果仍有争议
（因为技能本身也可能与社会经济背景相关），但它至少代表了向
公平招聘迈进的一步。

% -----------------------------------------------------------
\subsection{Moonhub：AI猎头的新范式}
\label{subsec:moonhub}
% -----------------------------------------------------------

Moonhub代表了AI招聘赛道的最新思路：\textbf{AI猎头}。与传统的
ATS（Applicant Tracking System）不同，Moonhub的AI不是在收到
简历后进行筛选，而是\textbf{主动在互联网上搜索和识别}潜在候选人
——分析LinkedIn档案、GitHub贡献、学术论文、专利申请、会议演讲等
公开数据，然后按照客户的需求进行排名和推荐。

Moonhub由斯坦福大学AI博士Nancy Xu创立，在硅谷的技术招聘圈中获得了
良好口碑。其模式更接近于"AI驱动的Executive Search"——适用于
高级管理层和稀缺技术人才的搜索，而非大规模的前线招聘（后者是
Paradox的领地）。

Moonhub的出现反映了AI招聘赛道的一个重要分化趋势——\textbf{高端
vs. 大众}。对于招聘一个年薪30万美元的高级工程总监，企业愿意花费
数万美元的猎头费用，也需要在全网范围内进行主动搜索——这是Moonhub
的价值区间。对于每天招聘50名仓库工人的物流公司，需要的是Paradox/
Olivia那种高效的漏斗自动化。对于评估已进入面试阶段的候选人，
需要的是HireVue的结构化评估。对于理解整个组织的技能分布和缺口，
需要的是Eightfold的智能分析平台。

每一层需求都足够大，足够支撑独立的产品线。但Workday、SAP等HR平台
巨头的战略是\textbf{全部整合在一个平台上}——Paradox的收购就是
这一整合策略的开端。

\subsection{AI HR赛道的公平性挑战}
\label{subsec:hr-fairness}

AI HR赛道面临着独特而严峻的公平性挑战——因为HR决策直接影响人的
生计。

\textbf{算法偏见}是最核心的问题。AI招聘系统是在历史数据上训练的，
而历史数据反映了过去的招聘偏见——如果一家公司过去90\%的工程师都是
男性，AI就可能学会偏好男性候选人。Amazon在2018年就因为发现其AI
招聘工具系统性地歧视女性候选人而废弃了该系统。

\textbf{可解释性}是第二个挑战。当AI拒绝一个候选人时，它需要能够
给出\textbf{可解释的理由}。"算法综合评分低于阈值"不是一个可接受
的理由——候选人有权知道自己被拒绝的具体原因。纽约市在2023年通过了
Local Law 144，要求使用AI招聘工具的雇主必须进行年度偏见审计并
公开结果。欧盟AI法案将AI招聘系统归类为"高风险AI"，要求更严格的
透明度和人类监督。

\textbf{数据隐私}是第三个挑战。AI招聘工具收集和分析大量个人数据
——简历、社交媒体档案、面试视频、甚至打字模式和鼠标移动。这些数据
的收集、存储和使用需要符合GDPR等隐私法规。特别是当Moonhub等工具
主动从公开来源抓取候选人信息时，是否需要候选人的明确同意，在不同
司法管辖区有不同的法律解释。

对于AI HR创业公司来说，公平性不仅是道德要求，也是\textbf{商业
必需}——一旦被证明存在系统性偏见，不仅会面临诉讼，更会失去
企业客户的信任。在这个赛道，"负责任的AI"不是营销口号，而是
生存前提。

\begin{keybox}[AI HR 与招聘赛道全景]
\small
\begin{tabularx}{\textwidth}{l c r X}
\toprule
\rowcolor{figLightGray}
\textbf{公司} & \textbf{成立} & \textbf{融资/估值}
  & \textbf{核心定位} \\
\midrule
HireVue     & 2004 & 私有
  & AI结构化面试; 从面部分析转型; 7000万+面试 \\
Paradox     & 2016 & 被Workday收购（2025/8）
  & 对话式AI招聘代理"Olivia"; 高流量前线 \\
Eightfold   & 2016 & \$397M / \$2.1B
  & 技能智能平台; D\&I; 三维度覆盖 \\
Moonhub     & 2023 & 早期融资
  & AI猎头; 主动搜索; 高端人才; 斯坦福血统 \\
\bottomrule
\end{tabularx}

\medskip
\noindent AI HR赛道的一个显著趋势是\textbf{整合加速}：Paradox
被Workday收购，Lattice（AI绩效管理）面临整合压力，iCIMS和
Greenhouse等ATS也在加速AI功能开发。独立AI HR创业公司的窗口期
可能只有2--3年——如果不能快速建立规模和客户粘性，就会被Workday、
SAP、Oracle等HR平台巨头吸收。对投资者来说，AI HR赛道的退出路径
越来越清晰：不是IPO，而是被大厂收购。
\end{keybox}

\bigskip

% =============================================================
% Chapter closing analysis
% =============================================================

回顾本章讨论的所有赛道——通用助手、AI伴侣、角色扮演、AI搜索、
AI浏览器、OS级AI、销售、客服、HR——一个统一的主题浮现出来：
\textbf{AI正在从工具变成入口，从应用变成基础设施}。

ChatGPT不再只是一个聊天机器人，它是一个平台——拥有Plugin生态、
GPT Store、桌面客户端和语音模式。Perplexity不再只是一个搜索引擎，
它是一个信息操作系统——从搜索到发布到任务执行的完整闭环。Apple
Intelligence不再只是一个功能，它是操作系统的一部分——覆盖超过
10亿台设备。Sierra的AI代理不再只是回答问题，它在\textbf{自主执行}
业务流程——处理退款、修改订单、升级服务方案。

\begin{corebox}[本章九大赛道的结构性洞察]
\begin{enumerate}
  \item \textbf{入口之争是终极之争}。苹果、Google、微软和字节跳动
    正在将AI嵌入操作系统和超级App，独立AI应用面临被截流的长期风险。
    ChatGPT桌面客户端和Perplexity IPO计划，都是在争夺独立入口。
  \item \textbf{用户规模 $\neq$ 商业价值}。豆包2.27亿月活但几乎
    不产生收入；Claude 1890万月活但年化收入140亿美元。高价值用户
    的密度决定了商业天花板。
  \item \textbf{AI伴侣是真实需求，但商业模式未解}。108亿美元的
    市场，39\%的年增长率，1.5亿用户——数字很漂亮。但Character.ai
    的3220万美元年收入和60\%的估值暴跌说明了一切：\textbf{参与度
    不等于付费意愿}。
  \item \textbf{AI搜索正在颠覆Google，但还没找到替代Google的商业
    模式}。Perplexity的200亿美元估值建立在对"搜索广告2.0"的想象上，
    但这个想象还需要数据验证。
  \item \textbf{AI客服是最接近"取代人类工作"的赛道}。Sierra和
    Intercom都预测2026年AI处理的客服量将超过人类——如果属实，
    这将是AI领域第一个大规模的人力替代事件。
  \item \textbf{Acqui-hire是AI创业公司的常见结局}。Inflection
    被微软掏空、Character.ai创始人被Google回收、The Browser Company
    被Atlassian收购、Paradox被Workday收购、Forethought被Zendesk
    收购——在大厂的资本和人才引力下，独立生存需要极强的战略定力。
  \item \textbf{中国市场的"规模大、收入小"困境需要解决}。豆包、
    DeepSeek、元宝合计月活近4亿，但变现远落后于ChatGPT和Claude。
    2026年春节的"红包大战"能否带来结构性改变，值得追踪。
  \item \textbf{开源正在重塑AI伴侣和角色扮演赛道}。SillyTavern的
    本地化方案和JanitorAI的无审查模式，正在吸纳从主流平台溢出的
    需求。随着本地模型质量接近云端，在线平台的付费价值被持续侵蚀。
  \item \textbf{"按结果付费"是SaaS商业模式的下一次革命}。从
    Sierra的"按解决量"到Intercom的"\$0.99/次"，AI客服赛道正在
    探索从"为人力付费"到"为结果付费"的根本转变。
\end{enumerate}
\end{corebox}

在这个演进过程中，最有价值的不是某一个AI产品的功能优势，而是
\textbf{谁拥有用户的习惯入口}。这就是为什么苹果、Google和微软
如此急切地将AI嵌入操作系统，为什么OpenAI要做桌面客户端和语音模式，
为什么Perplexity在考虑IPO，为什么Atlassian收购了The Browser
Company——入口之争，本质上是一场关于未来十年数字生活基础设施的
控制权之战。

这场战争有三种可能的结局：

\textbf{情景一："平台吞噬一切"}。苹果、Google和微软的OS级AI全面
成熟，独立AI应用被边缘化。用户不再需要打开ChatGPT或Perplexity，
因为Siri/Gemini/Copilot已经"足够好"了。在这个情景下，AI应用层
创业公司的价值被大幅压缩——就像独立杀毒软件被操作系统内置安全功能
边缘化一样。这对Perplexity和Character.ai是最坏的情景，但对Sierra
和Gong影响较小（因为企业级垂直场景的深度难以被通用平台替代）。

\textbf{情景二："AI版Google vs. Facebook"}。少数几家AI超级应用
（ChatGPT、Perplexity）建立了用户习惯和数据飞轮，与OS级AI形成
共存而非替代的关系。就像Google和Facebook在移动时代都成功独立于
iOS/Android生存一样，强大到足以抵抗平台的引力。在这个情景下，
AI应用层会产生2--3家千亿美元级别的独立公司。

\textbf{情景三："百花齐放的垂直化"}。通用AI助手趋同（谁都差不多好），
价值转移到垂直场景——AI法律助手、AI医疗助手、AI金融顾问、AI HR代理
各自建立了深度的行业壁垒。在这个情景下，不会有一家"AI版Google"，
但会有数十家在各自垂直领域做到数十亿美元收入的AI公司。这对Sierra、
Gong、Clay等垂直玩家是最好的情景。

真实的未来可能是三种情景的叠加——不同赛道、不同地理区域会走向
不同的结局。但有一点是确定的：\textbf{AI交互入口的争夺战才刚刚
开始}，而最终的版图将由用户的习惯——而非技术的先进性——来决定。

本章覆盖的九个赛道虽然跨度很大（从个人伴侣到企业客服），但它们
共享一条底层的进化轨迹：

\textbf{第一阶段（2022--2024）："助手"阶段}。AI是一个工具，你
向它提问，它给你答案。ChatGPT、Perplexity、Gong都始于这一阶段。
交互模式是单向的、被动的——人发起，AI响应。

\textbf{第二阶段（2024--2026）："代理"阶段}。AI开始主动执行任务。
Sierra的AI不只回答客户问题，还处理退款。Clay的AI不只富集数据，
还自动发送个性化邮件。ChatGPT的Operator不只搜索信息，还帮你
在网站上完成操作。这是当前正在发生的转变。

\textbf{第三阶段（2026+）："伙伴"阶段}。AI成为持续存在的智能层，
理解你的偏好、记住你的历史、预判你的需求。Apple Intelligence在
系统级别的集成、Nomi的成长型AI伴侣、以及Gong从"分析通话"到
"主动推荐下一步行动"的演进，都是这一阶段的先兆。

在这条进化轨迹上，每一步都在加深AI与用户之间的信任关系——而信任
一旦建立，就变成了最深的护城河。这也是为什么AI伴侣赛道虽然当前
收入微薄，但长期潜力可能被低估：它们在建立的不是工具依赖，而是
\textbf{情感依赖}——这种依赖的粘性远高于任何功能性工具。

对于正在阅读本章的创业者和投资者，我们的核心建议是：

\begin{enumerate}
  \item \textbf{不要在通用助手赛道正面硬刚大厂}。ChatGPT+Gemini+
    Copilot+豆包的组合已经覆盖了绝大部分通用场景。你的机会在垂直
    深度，而非通用广度。
  \item \textbf{关注"按结果付费"的商业模式创新}。Sierra和Intercom
    正在探索的定价模式可能代表了整个SaaS行业的未来方向。
  \item \textbf{AI伴侣赛道需要等待商业模式创新}。用户需求真实，
    但付费路径未通。下一个Character.ai级别的爆款可能不来自更好的
    模型，而来自更好的变现设计。
  \item \textbf{AI搜索的最大变量是Google的反应速度}。Perplexity的
    窗口期取决于Google AI Overviews的进化速度——如果Google足够快，
    窗口期可能只有2--3年。
  \item \textbf{OS级AI不会消灭独立应用，但会重新定义"默认值"}。
    当系统级AI覆盖了80\%的通用需求时，独立应用需要在剩下的20\%中
    做到10倍好。
\end{enumerate}

而对于这场战争中的创业公司来说，命运取决于一个简单的问题：
\textbf{你是在做一个会被大厂复制的功能，还是在建立一个他们
无法替代的习惯？}

最后，让我们回到本章开头提出的核心观察：这些产品覆盖了人与AI交互
\textbf{最高频、最私密、最具黏性}的场景。正是因为这些特征，本章
讨论的赛道比AI办公工具或AI创作工具更具战略意义——它们不只是在
改变人们\textbf{做事}的方式，而是在改变人们\textbf{获取信息}的
方式（搜索）、\textbf{建立关系}的方式（伴侣）、\textbf{与设备
交互}的方式（OS级AI）。

这些变化的长期影响可能远超当前的想象。当一整代人从青少年时期就
习惯于和AI角色对话、用AI搜索代替Google、在每一次手机交互中都有
AI参与时，他们对"信息""关系""隐私"的理解将会从根本上不同于
前一代人。本章覆盖的产品和公司，正站在这场认知革命的最前沿
——它们既是这场变革的推动者，也将是这场变革最大的受益者或牺牲品。

正如本书在第一章中所论述的，AI创业浪潮的本质是一次\textbf{基础设施
级别的范式转换}。本章讨论的"伴侣、搜索与个人助手"，正是这次转换
在消费者和前台业务层面最直接、最广泛的体现——从ChatGPT的9亿周活
到Character.ai的两小时平均使用时长，从Perplexity的7.8亿月查询
到Sierra的100亿美元估值，每一个数字都在讲述同一个故事：
\textbf{AI不再是未来，它已经是现在}。

下一章（第十章）将把视角转向企业端，探讨AI如何渗透到每一个行业的
每一个环节——从企业搜索到供应链优化，从网络安全到制造业质检。
如果说本章讨论的是AI如何改变"人与信息"和"人与人"的交互方式，
下一章将讨论AI如何改变"企业与效率"的关系。
